\section{Levels, boundary ends, and projective comparison}
\label{sec:D}

We use $\Gamma_{g,n}$ for the orientation-preserving mapping class group fixing every labelled puncture, $d=3g-3+n$, and $\mathcal C_{g,n}$ for the unmodified disjointness curve complex. For $0\leq k<d$, $P_k$ is its face poset restricted to multicurves of cardinality at least $d-k$. For $k\geq d$ we include the empty multicurve. This convention is essential at complexity one. The geometric realization of the empty-multicurve extension is a cone.

\subsection{The actual end and its fundamental group}

\begin{result}{D.1}{Proved}{The end comparison, including the exceptional types}
Let $2g-2+n>0$. Let $L_{g,n}$ be the boundary of the real oriented blow-up of the analytic Deligne--Mumford stack $\overline{\mathcal M}_{g,n}$, with corners rounded. The natural map
\[
 L_{g,n}\longrightarrow\mathcal M_{g,n}^{\rm an}\simeq B\Gamma_{g,n}
\]
has homotopy fibre $|\mathcal C_{g,n}|$. For nonempty boundary it induces an equivalence of fundamental groupoids precisely in the following ranges:
\[
 g=0,\ n\geq6;\qquad g=1,\ n\geq3;\qquad g\geq2,\ n\geq0.
\]
For $(0,5)$ and $(1,2)$ it induces a surjection on fundamental groups with kernel a free group of countably infinite rank. For $(0,4)$ and $(1,1)$ it does not induce a surjection on the fundamental group at a fixed boundary component. At $(0,3)$ the boundary is empty.
\end{result}
\begin{proof}
Choose an invariant thin part of Teichm\"uller space: a positive length threshold sufficiently small that curves shorter than the threshold are pairwise disjoint. The collar lemma gives this choice. For a nonempty multicurve $\sigma$, the locus where all its curves are short is contractible by Fenchel--Nielsen length and twist coordinates; intersections are the corresponding loci for unions of compatible multicurves, and are empty otherwise. The nerve of this cover is $\mathcal C_{g,n}$. A partition of unity provides the equivariant nerve comparison; equivariance can alternatively be obtained from the homotopy colimit of the cover, avoiding a choice of invariant partition on an infinite complex. The thin part is an equivariant model for the boundary collar of Harvey's bordification. Consequently the stack end is
\[
 E\Gamma_{g,n}\times_{\Gamma_{g,n}}|\mathcal C_{g,n}|,
\]
and its projection to $B\Gamma_{g,n}$ is the claimed homotopy-fibre sequence. This is the comparison used in EXIT\_ONE\_TYPE P5b1; it retains the stabilizer classifying spaces.

The exact sphere dimensions are
\[
 h(g,n)=\begin{cases}n-4&g=0,\\2g-3+n&g\geq1,\ n\geq1,\\2g-2&g\geq2,\ n=0.\end{cases}
\]
Harer's sphere theorem, in the explicit formulation of Birman--Broaddus--Menasco, Known Result~5 and equation~(1) \cite{BBM2015}, identifies the curve complex with a countably infinite wedge of $h(g,n)$-spheres. We have opened that formulation; Harer's original subscription text was not available. For $h\geq2$, the long exact sequence of the displayed fibre sequence gives connectedness and the asserted isomorphism of fundamental groups. For $h=1$, it gives
\[
 1\longrightarrow F_{\aleph_0}\longrightarrow\pi_1(L_{g,n})
 \longrightarrow\Gamma_{g,n}\longrightarrow1,
\]
since $\pi_2(B\Gamma_{g,n})=0$. For $h=0$, choose a curve $c$. Its component has group $\operatorname{Stab}_{\Gamma_{g,n}}(c)$, and the map is the stabilizer inclusion. The orbit of $c$ is infinite, so this subgroup is proper. These cases are exactly $(0,4)$ and $(1,1)$. At $(0,3)$ there are no stable graphs with an edge. Substituting the displayed dimensions proves the complete list.
\end{proof}

\begin{result}{D.2}{Proved}{Depth three is the sharp uniform input for the end's one-type}
In the normal-crossings link construction, retain every deployed stratum of depth at most $r$ and every flag among them. The resulting homotopy colimit maps naturally to the full link by an $(r-1)$-connected map. Thus depths $1,2,3$ compute its fundamental groupoid. Depths $1,2$ do not suffice uniformly: at the transverse meeting of three branches the fundamental groups are $\mathbb Z^4$ for the truncated diagram and $\mathbb Z^3$ for the full link. Applied to the stable moduli stack, this computes $\Pi_1(L_{g,n})$; D.1 gives precisely the types where it computes $\Pi_1(\mathcal M_{g,n})$. It does not compute $\Pi_1(\overline{\mathcal M}_{g,n})$ by that identification.
\end{result}
\begin{proof}
At a point on $q$ local branches, the normal angular factor is $(S^1)^q$. The auxiliary flag direction for the truncated diagram is the order complex of nonempty subsets of $\{1,\ldots,q\}$ of size at most $r$. This is the barycentric subdivision of the $(r-1)$-skeleton of a $(q-1)$-simplex. Its inclusion in the simplex is $(r-1)$-connected: the omitted cells have dimension at least $r$, and the skeleton is $(r-2)$-connected. Taking its product with $(S^1)^q$ gives the local comparison, including its homotopy fibre. The local contractions and their face maps are natural under permutations of the branches. They therefore descend through the finite stabilizer charts of the moduli stack and glue in the deployed flag homotopy colimit. Relative cellular connectivity is preserved under these gluings. For $q=3,r=2$ the auxiliary complex is the boundary of a triangle. Thus the local truncated model is $(S^1)^3\times S^1$, whereas the full auxiliary simplex is contractible. The displayed groups follow by the product formula. The map kills the last factor and is the identity on the three normal meridians. This proves both sharpness and the exact corrected target of cote~147, scan~152, formula~(247).
\end{proof}

\begin{result}{D.3}{Proved}{Compactification independence of the two end constructions}
For any compact Hausdorff compactification $U\subset X$ with $D=X\setminus U$, the archimedean punctured-neighborhood pro-space is canonically the end system $\{U\setminus K\}_{K\subset U\text{ compact}}$. For smooth proper normal-crossings pairs this pro-homotopy type is constant, with representative the rounded real-oriented-blow-up boundary. Over a trivially valued field $k$, a proper modification $X'\to X$ equal to the identity on $U$ induces an isomorphism
\[
 W(X',X'\setminus U)\xrightarrow{\ \sim\ }W(X,D),\qquad
 W(X,D)=X^{\beth}\setminus(U^{\beth}\cup D^{\beth}).
\]
For proper $X$, this space is $U^{\rm an}\setminus U^{\beth}$. Properness cannot be removed: the compactifications $\mathbb G_m\subset\mathbb A^1$ and $\mathbb G_m\subset\mathbb P^1$ give one and two punctured boundary tubes respectively.
\end{result}
\begin{proof}
The mutually inverse assignments are $N\mapsto K=X\setminus N$ and $K\mapsto N=X\setminus K$, where $N$ is an open neighborhood of $D$. Compact Hausdorffness makes each $K$ compact and each embedded compact $K$ closed. They identify the actual punctured open sets and every refinement map. In the normal-crossings chart $z_1\cdots z_q=0$, the oriented blow-up has coordinates $z_i=\rho_i e^{i\theta_i}$ with $\rho_i\geq0$. A rounded boundary collar gives the cofinal end family $L\times(0,\epsilon)$ and all bonding maps are homotopy equivalences.

For the nonarchimedean assertion, a point of $X^{\beth}$ is a seminorm with a center on $X$. A point of $W(X,D)$ has its generic point in $U$ and its center in $D$. Properness gives a unique lift of that center to $X'$, and separatedness makes the lift unique. Both tubes are thereby identified with the same analytic domain of $U^{\rm an}$, including its induced analytic structure. This is the valuative proof of Ben-Bassat--Temkin's comparison; see \cite[\S3.3.2 and Proposition~4.5.1]{BBT2013}. For $\mathbb G_m$ the two explicit domains are $0<|t|<1$ and the disjoint union of that domain with $1<|t|<\infty$. The natural map of the first to the second is the inclusion of one connected component. These comparisons prove the corrections to the added claim about cote~149, scan~133, and verify LM5 F.40--41 in the stated meanings.
\end{proof}

\begin{result}{D.4}{Proved}{The two requested ends and their comparison maps}
For the tripod, the archimedean end is three circles, the analytic open curve has homotopy type $S^1\vee S^1$, and the underlying trivially valued punctured tube is three open intervals. For $(\overline M_{0,5},D)$ the integral homology is free, with Betti numbers
\[
\begin{array}{c|rrrr}
 &b_0&b_1&b_2&b_3\\\hline
 M_{0,5}^{\rm an}&1&5&6&0\\
 L_{0,5}&1&11&11&1\\
 W(\overline M_{0,5},D)&1&6&0&0.
\end{array}
\]
The first comparison in D.1 sends every peripheral circle to its meridian. The nonarchimedean skeleton map sends a valuation to the normalized vector of its boundary orders, and retracts its underlying space onto the Petersen graph.
\end{result}
\begin{proof}
The tripod calculation follows by taking three punctured disks and a pair-of-pants core. A seminorm on a punctured formal disk over a trivially valued algebraically closed field is determined by $r=|t|\in(0,1)$: every unit has norm one, and a series of order $m$ has norm $r^m$.

Normalize three labels of $M_{0,5}$ to $0,1,\infty$. Projection $(x,y)\mapsto x$ from $\{x,y\in\mathbb C\setminus\{0,1\}:x\ne y\}$ is a bundle with base a twice-punctured plane and fibre a thrice-punctured plane. The monodromy fixes the oriented puncture meridians in first homology. The integral Serre spectral sequence has only two rows and columns and no possible differential; its ranks are the coefficients of $(1+2t)(1+3t)$.

The ten boundary divisors are indexed by two-element subsets of five labels; two meet exactly when the subsets are disjoint. They meet transversely once, and no three meet. Each is a rational curve of self-intersection $-1$ in the degree-five del Pezzo model. A regular neighborhood $N$ retracts to this union, so $H_1(N)=\mathbb Z^6$, $H_2(N)=\mathbb Z^{10}$, and $H_3(N)=0$. The intersection matrix is $Q=-I+A$, where $A$ is the Petersen adjacency matrix. The exact integral computation in \texttt{code/d\_boundary\_checks.py} gives Smith form with five unit diagonal entries and five zero entries. Poincar\'e--Lefschetz duality and the exact sequence of $(N,\partial N)$ then give $H_1(L)=\mathbb Z^{11}$ and $H_2(L)=\mathbb Z^{11}$; connectedness and orientation give $H_0=H_3=\mathbb Z$. This independently verifies LM5 F.42's archimedean row.

The local skeletal coordinates of the nonarchimedean tube are $a_i=-\log|z_i|\geq0$, not all zero. Normalization by $\sum a_i$ defines the comparison to the boundary complex. The toroidal retractions glue under boundary-coordinate changes and give the deformation retraction proved in Thuillier's theorem \cite[Theorem~3.25 and Proposition~4.5]{Thuillier2007}. The Petersen graph has ten vertices, fifteen edges and first Betti number six, proving the last row.
\end{proof}

\subsection{Rank selection and the levels}

\begin{result}{D.5}{Proved}{The exact link calculation and a uniform rank-selection bound}
For a nonempty multicurve $\sigma$ let $v_0(\sigma)$ be the number of rational components after cutting. Then its upper link has homotopy type a wedge of spheres of dimension
\[
 H(\sigma)=2g-3+n-v_0(\sigma).
\]
A pair-of-pants factor is interpreted as the join unit $S^{-1}$. For $g\geq1$ and $0\leq k<d$, $|P_k|$ is at least
\[
 \left(\min\left\{h(g,n),\left\lfloor\frac{2(k+1)}3\right\rfloor\right\}-1\right)\text{-connected}.
\]
For $g=0$ it is $(k-1)$-connected. In every genus $P_1$ is connected and $P_2$ is simply connected, with the empty-multicurve convention at $k\geq d$.
\end{result}
\begin{proof}
Write the cut components as $S_{g_i,n_i}$. Every component has $n_i>0$ and their complexes join to the link. The sphere dimension of their join plus one is the sum of the individual sphere dimensions plus one. With $v_+$ positive-genus vertices, $v=v_0+v_+$, and $q=|\sigma|$, the identities
\[
 \sum n_i=n+2q,\qquad \sum g_i=g-q+v-1
\]
give
\[
 \sum_{g_i>0}(2g_i-2+n_i)+\sum_{g_i=0}(n_i-3)-1
 =2g-3+n-v_0.
\]
This also directly types the join comparison; it is obtained by restricting disjoint curves to the cut components and taking their union for the inverse map.

Let $r=\sum(3g_i-3+n_i)=d-q$ be the residual complexity. A positive-genus factor of complexity $r_i>0$ has $g_i\leq\lfloor(r_i+2)/3\rfloor$, hence sphere dimension plus one at least $\lfloor2r_i/3\rfloor+1$. A rational factor contributes exactly $r_i$. For each positive integer $r_i$ these quantities are at least $2r_i/3+1/3$. Summing and taking the integer part proves $H(\sigma)\geq\lfloor2r/3\rfloor$ whenever $r>0$; this also holds when all factors are rational, since then $H=r-1\geq\lfloor2r/3\rfloor$ for $r\geq3$. The only smaller values $r=1,2$ are checked directly and give $H=0,1$.

Delete barycentric vertices of sizes $q=1,2,\ldots,d-k-1$ in that order. At deletion the remaining link is the upper link above and $r\geq k+1$. For $k\geq2$ the lower bound just proved applies; for $k=0,1$ use nonemptiness and the direct small-rank checks. Coning a $c$-connected link is a $(c+1)$-connected attachment and preserves homotopy groups through degree $c$. The initial curve complex is $(h-1)$-connected. Vertices of equal rank are incomparable, so all cones at one rank attach independently; weak CW realizations justify the infinite coproduct. This yields the asserted estimate. In genus zero every link has dimension $d-q-1$ and connectivity $d-q-2$, giving the sharper $(k-1)$ bound. For $k=1,2$ the displayed estimate gives respectively $0,1$; the exceptional low complexities are cones.
\end{proof}

\begin{result}{D.6}{Refuted}{The original Lead 4 and the four requested computations}
The original bound $(\min(k,2g-2+n)-1)$ is false: at $(g,n,k)=(2,1,3)$ the complex is a nontrivial wedge of two-spheres. The complete low-type table, with $\ast$ denoting a contractible cone and $\bigvee_\infty$ a countable wedge, is
\[
\begin{array}{c|c|c|c|c}
(g,n)&d&P_1&P_2&P_3\\\hline
(1,2)&2&\bigvee_\infty S^1&\ast&\ast\\
(1,3)&3&\bigvee_\infty S^1&\bigvee_\infty S^2&\ast\\
(2,0)&3&\bigvee_\infty S^1&\bigvee_\infty S^2&\ast\\
(2,1)&4&\bigvee_\infty S^1&\bigvee_\infty S^2&\bigvee_\infty S^2.
\end{array}
\]
All $P_0$ are nonempty discrete sets; $P_k$ is contractible for $k\geq d$. The corrected proposed bound with $h(g,n)$ in place of $2g-2+n$ holds in all four rows.
\end{result}
\begin{proof}
At $k=d-1$ the poset is all nonempty faces of the curve complex, so barycentric realization and the opened sphere theorem give the wedge entries. For $k<d-1$, the dimension of $P_k$ is $k$ and D.5 gives connectedness at $k=1$ and simple connectedness at $k=2$. Each connected countable graph is a wedge of circles. For $(1,3)$ and $(2,0)$ cut one nonseparating curve; its complement is respectively $(0,5)$ and $(1,2)$. The proper cofaces of that curve embed the barycentric subdivision of its curve complex as a subgraph of $P_1$. For $(2,1)$ use two nonseparating curves with connected complement $(0,5)$. These subgraphs have countably infinite first Betti number, and inclusion of a subgraph injects first homology, proving all circle multiplicities. For $P_2$ at $(2,1)$ cut one nonseparating curve. Its complement has type $(1,3)$, so its proper-coface subcomplex has countably infinite second homology. Inclusion into the two-dimensional $P_2$ injects second homology, because no three-chains exist. A simply connected two-dimensional CW complex is a wedge of two-spheres: choose a basis of its free second cellular homology, realize it by spheres using Hurewicz, and apply the homological Whitehead theorem. Countability then gives exactly the displayed multiplicity. At $(2,1,3)$ the second homology is a countably generated free abelian group, so the original asserted two-connectedness fails. This is an infinite topological counterexample; the finite calculations supplied with this section verify the graph types and rank data, not a replacement of the infinite curve complex by its finite quotient.
\end{proof}

\begin{result}{D.7}{Open}{The stronger all-genus rank-selected connectivity claim}
For arbitrary stable $(g,n)$, the claim that $P_k$ is $(\min(k,h(g,n))-1)$-connected has not been established here. D.5 proves a uniform lower bound and D.6 settles every requested low type. The first obstruction to the proposed proof occurs at residual complexity four: a component of type $(2,1)$ has a link with nonzero second homology, so deleting its associated barycentric vertex is only two-connected. For example, in type $(3,0)$, a two-curve cut can leave a $(2,1)$ component and a pair of pants. The proposed $k=3$ bound there is two-connectedness, while the direct link estimate supplies only simple connectedness.
\end{result}

We tested the exact upper-link formula, the uniform rank-deletion argument, and finite stable-graph enumerations. The finite quotient does not decide the missing absolute homotopy group: the quotient map retains mapping-class stabilizers, and its effect on $\pi_2$ has no asserted injectivity. The next necessary calculation is the boundary map from the homology created by the preceding rank-one deletions to the sums of the $(2,1)$ upper-link Steinberg modules. No surjectivity of this map is used in the proved statements.

\subsection{Checking the projective threefold comparison}

In this subsection a \emph{good pair} is a smooth complex projective variety with reduced normal-crossings boundary, every open boundary stratum aspherical, and every full normal-meridian lattice injecting into $\pi_1(U)$. The marking by $U$ is part of the data. The comparison below concerns these algebraically marked completions. It makes no assertion about arbitrary birational maps that move $U$.

\begin{result}{D.8}{Proved}{Good projective pairs are Mori hyperbolic and their comparisons are log crepant}
A good projective pair has no nonconstant algebraic map $\mathbb A^1\to S$ into any open stratum. For a simple normal-crossings boundary, $K_X+D$ is nef. For two such completions and a common smooth normal-crossings resolution $p:Z\to X$, $q:Z\to Y$, write $B=Z\setminus U$. Then
\[
 K_Z+B=p^*(K_X+D)+A_p=q^*(K_Y+E)+A_q,
 \qquad A_p=A_q\geq0.
\]
The common divisor is exceptional for both projections. A proper birational morphism between good smooth projective simple normal-crossings pairs, equal to the identity on $U$, is toroidal and log \'etale.
\end{result}
\begin{proof}
A map $f:\mathbb A^1\to S$ extends to $\bar f:\mathbb P^1\to X$ by properness. If infinity remains in $S$, asphericity makes the sphere nullhomotopic, whereas the integral of the pullback of a projective K\"ahler form is positive for a nonconstant holomorphic sphere. If infinity enters a deeper stratum $T$, trivialize the pulled-back old normal lines on the contractible source $\mathbb A^1$. Push the map a small nonzero distance in these directions into $U$, choosing the radius to tend to zero faster than every newly vanishing boundary coordinate near infinity. The boundary of a large disk then gives a nullhomotopic loop in $U$. In a normal-crossings chart at $\bar f(\infty)$ its newly acquired meridian coordinates are the nonzero contact orders of $\bar f$, with a common orientation sign. The old normal coordinates do not alter those new coordinates. This contradicts injection of the full lattice at $T$. Thus no such map exists. The open strata are exactly the strata in Svaldi's Theorem~1.1 \cite{Svaldi2019}; its hypotheses are smoothness, projectivity, reduced SNC boundary and the just-proved absence of these affine lines. Its conclusion gives nefness.

The logarithmic differential maps have effective determinant divisors $A_p,A_q$ supported on exceptional prime divisors. Put
\[
 R=p^*(K_X+D)-q^*(K_Y+E)=A_q-A_p.
\]
The divisor $-R$ is $p$-nef and $p_*R=p_*A_q$ is effective. The negativity lemma therefore gives $R\geq0$. Interchanging $p$ and $q$ gives $-R\geq0$, hence equality. This proves the claimed identity without asserting that the common resolution is good.

For an actual birational morphism $f:Y\to X$, the effective logarithmic Jacobian $A$ in $K_Y+E=f^*(K_X+D)+A$ is exceptional and $f$-nef. The same negativity lemma gives $A=0$. Thus the logarithmic differential matrix is everywhere invertible. In local boundary coordinates, $f^*x_i=u_i\prod_j y_j^{a_{ij}}$. The exponent matrix has full column rank; after choosing a nonsingular square minor and taking \'etale roots of units, absorb the corresponding units into the $y_j$. The remaining boundary-unit functions together with the ordinary target coordinates have independent ordinary differentials. They provide the invertible coordinates of a toric chart. Hence the morphism is locally an \'etale pullback of a monomial toric morphism; the finite lattice cokernel is invertible in characteristic zero. This is the explicit chart proof of log \'etaleness and toroidality.
\end{proof}

\begin{result}{D.9}{Proved}{Theorem R for projective threefolds, with the factorization reference corrected}
Two good smooth projective normal-crossings completions $(X,D)$ and $(Y,E)$ of the same smooth complex threefold $U$ admit a common good smooth projective completion $X_s$ such that $X_s\to X$ is a sequence of blow-ups of smooth closures of strata and $X_s\to Y$ is toroidal. They are joined over $U$ by a finite sequence of smooth stratal blow-ups and blow-downs through good projective pairs. This confirms Theorem~R of the attached review in its projective threefold scope.
\end{result}
\begin{proof}
First blow up the closures of the self-intersection strata in increasing dimension, obtaining simple normal-crossings boundaries. Locally these are barycentric subdivisions of the normal orthants. Every new open stratum is a torus bundle over an old open stratum, hence is aspherical by the homotopy exact sequence. Its normal-meridian lattice maps by a linearly independent set of primitive fan vectors into the old lattice, hence still injects into $\pi_1(U)$. This proves goodness at every step, and also proves the same preservation for the further stratal blow-ups below.

By D.8 every component $E_j$ of $E$ has log discrepancy zero over $(X,D)$. For a smooth SNC pair, the divisorial valuations of log discrepancy zero are exactly those produced by successive blow-ups of strata. One way to see the equality criterion is to pull back the logarithmic volume form: a center lying properly inside an open stratum supplies at least one ordinary coordinate and gives strictly positive discrepancy. A sequence following the center of the valuation therefore has discrepancy zero only at corner centers. In a normal orthant this is precisely a primitive monomial valuation. This criterion is also stated after Definition~2.7 in Svaldi's opened text, by reduction to the SNC discrepancy calculation. A primitive integral ray can be extracted by smooth star subdivisions: in a chart containing the ray, subtract the least positive coordinate from the others, reducing their sum until the ray is a coordinate ray. Finitely many such extractions give a good smooth projective $X'\to X$ on which every prime divisor of $Y$ occurs.

Choose a very ample line bundle $H$ on $Y$. Transport its complete linear system to $X'$ and remove the common divisorial part. Let $I$ be the remaining base ideal and $L$ its line bundle. Its cosupport lies in $D'=X'\setminus U$ and has codimension at least two. Let
\[
 p:Z=\operatorname{Nor}(\operatorname{Bl}_I X')\longrightarrow X',
 \qquad q:Z\longrightarrow Y
\]
be the normalized graph. Writing $I\mathcal O_Z=\mathcal O_Z(-F)$, we have $q^*H=p^*L-F$ and $\mathcal O_Z(-F)$ is $p$-ample. Consequently $q$ is finite on every fibre of $p$: on such a fibre, the pullback of $H$ is ample, and a positive-dimensional fibre of $q$ would contain a curve on which that line bundle has degree zero.

Let $A$ be a $p$-exceptional prime divisor, equivalently a Rees divisor. It is absent as a divisor of $X'$, and therefore absent as a divisor of $Y$ by our extraction. Thus $\dim q(A)\leq1$. Its center on $X'$ cannot be a point, since then the two-dimensional $A$ would be mapped finitely by $q$ into a curve or a point. Hence its center is a curve $B$, and $q(A)$ is a curve $R$.

We now justify the crucial assertion that $B$ is a stratum curve and that the valuation is monomial. Choose a general point $b\in B$ avoiding every other exceptional center that does not dominate $B$, and a smooth surface slice transverse to $B$. Resolve the graph by principalizing $I$ on a smooth model. On the slice the resolution is a succession of point blow-ups of a smooth surface. Its exceptional divisor has a tree as dual graph, and its components are rational curves. The strict transforms of the boundary branches through $b$ give either one or two marked arrows. Each branch gives one arrow because its smooth curve germ has a single strict transform meeting the exceptional fibre in a single point. There are at most two branches at the generic point of a curve in a smooth SNC threefold.

Retain the vertices corresponding to Rees components, together with the marked arrows, and let $T_0$ be their minimal connecting subtree. A Rees component maps finitely to a rational curve $R_i\subset Y$. Let $S_i$ be the open stratum containing its generic point. The normalization of $R_i$ has at least two points outside $S_i$, by the absence of maps $\mathbb A^1\to S_i$ proved in D.8 (zero missing points also gives such a map). The same holds on its source rational component, because a finite surjection has a nonempty fibre over each of those two distinct points.

Every such boundary point comes from a boundary component $E_j$ not containing $R_i$. Its inverse image is an effective Cartier divisor on the normal graph $Z$, and therefore has pure codimension one. Here is the precise purity argument needed. Any supporting prime exceptional over $X'$ is a Rees divisor. Its center is a curve by the preceding paragraph; since $b$ avoids every other such center, that curve must be $B$. A supporting prime which is not exceptional is the strict transform of a prime divisor of $X'$. The map is the identity on $U$, so that prime lies in $D'$; it meets the general fibre over $b$ only when it contains $B$. These are exactly the one or two marked branches. The chosen Rees component is not itself in this support, since $E_j$ does not contain its image $R_i$.

On the resolved transverse surface, every exceptional component not corresponding to a Rees divisor contracts to a point of $Z$: a component mapping onto a curve of the general fibre would sweep out an exceptional prime of $Z$ dominating $B$, and would itself be a Rees component. Its connected non-Rees subtrees likewise contract. Thus all arrows and Rees components in any one direction of the resolved tree meet the chosen component at one point after those intervening subtrees are collapsed. A branch containing no arrow and no Rees component cannot supply an extra boundary point, because it supplies no supporting prime of the Cartier pullback on $Z$. It follows that every Rees vertex has at least two incident directions in $T_0$.

A finite tree with a retained vertex and only one marked arrow has an unmarked leaf. This contradicts the preceding assertion. Thus there are two arrows, so $B$ is an irreducible component of the intersection of two boundary divisors. Moreover $T_0$ is the path between the two arrows: a branch off this path would end at a Rees leaf. In the resolution of two transverse smooth curve germs, the vertices on that path are exactly those obtained by successive blow-ups of corner intersection points. Induction verifies this: a corner blow-up subdivides the path, while a blow-up at any other point adds a branch off it. Consequently every Rees valuation is a monomial valuation in the two boundary parameters.

For every $m\geq1$, the integral closure of $I^m$ is the intersection of the valuation ideals $v(f)\geq m v(I)$ over these finitely many Rees valuations. We use the existence and normalized-blow-up construction in Huneke--Swanson, Definition~10.1.1 and Theorem~10.2.2 \cite{HunekeSwanson2006}, both opened. Each inequality is monomial in boundary parameters, so each integral closure, and in particular $\overline I$, is a toroidal ideal sheaf. This includes points where several center curves meet. In a triple-crossing chart $x,y,z$, a Rees valuation over $x=y=0$ has weights $(a,b,0)$, and its valuation ideal is generated by the monomials $x^i y^j z^k$ with $ai+bj\geq m v(I)$. Valuations over $x=z=0$ or $y=z=0$ give the other two kinds of monomial inequality. At the special point the local integral closure is the intersection of all these valuation ideals. A finite intersection of monomial ideals is monomial: its generators are the least common multiples of choices of generators. The same argument applies to every subcollection of the stratum curves. An additional isolated integral-closure condition would require a Rees valuation centered at that point, already excluded by finiteness on $p$-fibres. Thus the normalized blow-up is toroidal everywhere, including the special triple points.

Principalize this monomial ideal by smooth stratal blow-ups. Canonical principalization, in the form recalled in AKMW \S1.2.3 \cite{AKMW2002}, is invariant under local torus actions and \'etale charts. Applied to a monomial ideal its smooth centers are unions of torus orbits closed in the chart, hence closures of strata; the local procedures agree on overlaps. We obtain a good $X_s\to X'$ and a morphism $X_s\to Y$. D.8 makes the latter toroidal. The corrected toroidal weak factorization theorem, AKMW Theorem~2.7.1 and Matsuki's erratum \S2 \cite{Matsuki2000}, factors it into smooth stratal blow-ups and blow-downs.

We use its relative morphism case. For the projective toroidal morphism $X_s\to Y$, the cobordism has a fan map to the fan of $Y$; its subdivisions preserve this map, and every elementary quotient fan still maps to that endpoint fan. These are the explicit fan maps in the proof of AKMW Theorem~2.7.1. Thus every intermediate variety is projective over $Y$ and its fan subdivides the fan of $Y$. This relative assertion proves goodness through the blow-down steps: every intermediate open stratum is a torus bundle over an open stratum of $Y$, and its meridian lattice injects, by the independent primitive generators of its cone, into the normal lattice of that endpoint stratum. Asphericity and meridian injectivity follow directly from goodness of $Y$, independently of the direction of the adjacent blow-up arrow. This gives the entire stated zigzag.

The proof uses only weak factorization. The algorithm for strong toroidal factorization in the earlier Morelli/AMR citation was withdrawn in the opened erratum and is not used.
\end{proof}

\begin{result}{D.10}{Proved}{All one-step stratal modifications of $\overline M_{0,6}$}
There are exactly $210$ nontrivial labelled smooth stratal blow-ups of $\overline M_{0,6}$: $105$ along boundary curves and $105$ at boundary points. Under permutation of the six labels they form four orbits, of sizes $45,60,15,90$. There is no nontrivial one-step stratal blow-down from the standard pair to another smooth SNC pair. The associated maps of conical complexes are the star subdivisions at the chosen two- or three-dimensional cones, with new ray the sum of their primitive generators. Theorem D.9 consequently applies to every good smooth projective completion of $M_{0,6}$.
\end{result}
\begin{proof}
A stable six-legged tree has at most three edges. The exact enumeration in \path{code/d_boundary_checks.py}, using the independently generated graph list of Part~A, gives $25,105,105$ strata of codimensions one, two, three. Blowing up a divisor is the identity. For two-edge trees the end/middle leg distributions are $(2,2,2)$ and $(2,1,3)$, with $45$ and $60$ labellings. For three-edge trees the internal graph is either a three-armed star, with two legs at each endpoint, or a path with endpoint leg counts two and interior leg counts one. Their counts are $15$ and $90$. Smoothness of the corresponding star subdivisions follows because each replacement maximal-cone matrix has determinant $1$; all these determinants are checked exactly by the program.

A smooth stratal blow-down in dimension three must contract a boundary divisor that is either a projective plane over a point or a projective-line bundle over a smooth curve. A divisor of type $2|4$ is $\overline M_{0,5}$, a rational surface with second Betti number five. It is neither $\mathbb P^2$ nor a projective-line bundle over a curve: rationality forces that curve to have genus zero, and such a bundle has second Betti number two. A divisor of type $3|3$ is $\mathbb P^1\times\mathbb P^1$. Any contraction to a curve uses a ruling, but its boundary contains three sections of that ruling. Their images would be three distinct boundary branches containing the entire center curve. Three SNC branches in a smooth threefold meet in codimension three, so this target would fail the SNC condition. These exhaust the twenty-five divisors. Finally, at each stratum of the standard moduli compactification the cut-surface Teichm\"uller groups are aspherical and independent commuting Dehn twists inject; hence the pair is good and D.9 applies.
\end{proof}

\subsection{Completion of the level complexes}

Fix a finite-index pure subgroup $G\subset\Gamma_{g,n}$ so that the simplicial action has no inversions; taking a barycentric subdivision also removes inversions. For a simplicial $G$-complex $K$ with finitely many cell orbits, its \emph{full level completion} is the simplicial profinite set
\[
 K^{\mathrm{lev}}=\varprojlim_{H\leq_fG}K/H.
\]
The procongruence version uses only congruence subgroups $H$. The map between them is induced by restriction of the inverse system. These definitions fix the exact completions used below.

\begin{result}{D.11}{Proved}{Profinite simple connectivity at level two and above}
For every stable type and every $k\geq2$, the full level completion of $P_k$ has trivial profinite fundamental group. The completed $P_1$ is connected. For the full curve complex, the same simple-connectivity conclusion holds for $d\geq3$. Completing the discrete presentation from the complex of groups on $P_2$ gives the full profinite group $\widehat\Gamma_{g,n}$ as a profinite presentation, with generators furnished by strata of modular dimension at most one and additional gluing relations furnished in dimension two.
\end{result}
\begin{proof}
D.5 and the cone convention show that $P_k$ is simply connected for $k\geq2$. For any finite-index $H\subset G$, Armstrong's quotient theorem gives
\[
 \pi_1(P_k/H)=H/H_*,
\]
where $H_*$ is the normal subgroup generated by elements with fixed points. The theorem's hypotheses hold: the complex is simply connected, the action is simplicial and without inversions. The same theorem and derivation are opened in Kent \S7, Theorem~11 \cite{Kent2016}. Thus $\widehat H$ surjects onto $\widehat{\pi_1(P_k/H)}$. Inverse limits of these surjections are surjective by compactness. But $\varprojlim_{H\leq_fG}\widehat H=1$: each $\widehat H$ identifies with the closure of $H$ in $\widehat G$, and their intersection is the identity by the defining finite-quotient topology. The profinite fundamental group commutes with this inverse limit of finite simplicial sets, in the sense specified in Kent \S5. Hence the desired group is trivial. The proof for the curve complex uses its simple connectivity at $d\geq3$. Connectedness of completed $P_1$ follows from connectedness of all its finite-level quotients and the corresponding inverse-limit statement for $\pi_0$.

For the presentation assertion, choose a maximal tree, stabilizer lifts, and transporter elements in the discrete complex of groups. Its ordinary fundamental group is $\Gamma_{g,n}$ because $P_2$ is simply connected. The homomorphisms from this presentation to any finite group are exactly the compatible homomorphisms from the local stabilizers and the edge transporters satisfying the edge and triangle relations. By the universal property of completion, homomorphisms from each local group to a finite group extend uniquely and continuously to its full completion. The same finite-group-valued universal property therefore identifies the resulting profinite presentation with $\widehat\Gamma_{g,n}$. This proves the presentation as a profinite colimit; it does not silently assert injectivity of every completed local-group map. Such injectivity is a separate comparison supplied where needed by Part~E.
\end{proof}

\begin{result}{D.12}{Refuted}{A printed low-complexity range in the profinite literature}
For the unmodified disjointness complex, the full level completion of $\mathcal C_{0,5}$ has a nontrivial continuous $\mathbb Z/2$-valued first cohomology class, and hence nontrivial profinite fundamental group. Thus the bound $3g-3+n\geq2$ printed in Kent's Theorem~12 cannot hold for the complex defined there. The argument in D.11 proves the corrected range $3g-3+n\geq3$. The genus-zero case of the analogous printed range in Theorem~13 likewise requires this correction.
\end{result}
\begin{proof}
Place the five labelled points on a circle. Let $c_{12},c_{23},c_{34},c_{45},c_{51}$ be boundaries of small regular neighborhoods of the five consecutive arcs. Curves associated with disjoint pairs are disjoint. Therefore
\[
 c_{12},c_{34},c_{51},c_{23},c_{45},c_{12}
\]
is a closed pentagon in the actual infinite curve complex. The quotient by the pure mapping class group maps a curve to the two-element side of its partition of the labels. It is the Petersen graph: edges join disjoint two-element subsets. The pentagon above maps to an embedded pentagon in that graph. On the Petersen graph take the $\mathbb Z/2$ cellular one-cochain equal to one on the edge $\{12,34\}$ and zero on all other edges. Every such cochain is a cocycle, since the graph has no two-cells, and its evaluation on this pentagon is one.

For every finite-index subgroup $H$ of the pure mapping class group, there are compatible maps
\[
 \mathcal C_{0,5}\longrightarrow\mathcal C_{0,5}/H
 \longrightarrow\mathcal C_{0,5}/\Gamma_{0,5}.
\]
The displayed closed pentagon supplies a closed lift in every intermediate graph. Thus the pulled-back cocycle remains nonzero at every level. It defines a nonzero continuous cocycle on the inverse-limit simplicial profinite graph: if it were a coboundary, a continuous zero-cochain with finite target would factor through one finite level, where the pentagon evaluation rules out a coboundary. This gives the asserted continuous character of the profinite fundamental group.

The finite Petersen graph, the lifted pair pattern, and the cocycle evaluation are checked in \texttt{code/d\_boundary\_checks.py}. No computation on a finite quotient is being used to prove that an arbitrary infinite loop survives: we supplied the actual closed lift. The source of the printed problem is visible in Kent's Theorem~10, which gives the correct Harer sphere dimensions but writes the complex dimension as $3g-3+n$ and infers simple connectivity one complexity too early. Armstrong's theorem requires that missing simple connectivity. The same sphere theorem gives $\mathcal C_{0,5}\simeq\bigvee_\infty S^1$ and D.11 supplies the corrected elementary argument.
\end{proof}

\begin{result}{D.13}{Known}{The proved procongruence range and the remaining scope}
For $n\geq2$, $2g-2+(n-1)>0$, and $3g-3+n\geq3$, the procongruence curve complex has trivial profinite fundamental group. This follows from Kent's Theorem~13 with its proof restricted to the actual simply connected range. The present proof does not establish the stronger higher-connectivity claims for arbitrary procongruence $P_k$, nor does it infer the congruence subgroup property from this one complex.
\end{result}
\begin{proof}
Put $S=S_{g,n-1}$ and $\dot S=S_{g,n}$. The first two inequalities make $S$ a smooth affine hyperbolic curve; the third makes $\mathcal C(\dot S)$ simply connected, which is the exact hypothesis required for the Armstrong sequence used in the proof. We spell out the comparison employed in Kent's \S8. Let $K$ be the congruence kernel, and let $\kappa$ range through finite-index congruence subgroups of the pure level subgroup. Armstrong's theorem, profinite right exactness, and inverse-limit exactness give
\[
 \varprojlim_\kappa\widehat{\kappa_*}\longrightarrow K
 \longrightarrow\Pi_1(\check{\mathcal C}(\dot S))\longrightarrow1.
\]
Choose a curve $\gamma$ bounding a pair of pants containing the added puncture and another puncture. Kent's Theorem~3 and Corollary~4 prove that the products $B_\kappa\kappa_\gamma$ are congruence and cofinal, where $B$ is the point-pushing subgroup and $B_\kappa=B\cap\kappa$. In the opened proof of Lemma~15, closure separation of the stabilizer implies
\[
 \bigcap_\kappa\overline{B_\kappa h_\gamma}
       =\overline{h_\gamma}
\]
for every fixed congruence subgroup $h$. Cofinality then places $K$ in the image of $\varprojlim\widehat{\kappa_\gamma}$, which lies in the image of $\varprojlim\widehat{\kappa_*}$. The displayed quotient is trivial. This is the exact step of the cited proof that transfers to our conventions. It uses the affine Birman sequence and the congruence cofinality proved in the source, not an assertion that all finite-index subgroups are congruence. D.12 explains why we have stated the necessary complexity restriction explicitly.
\end{proof}

\subsection{The twenty-three lambencies, computed as covers}

\begin{result}{D.14}{Proved}{All passports, cusp fields, Atkin--Lehner actions, and signatures}
For each lambency in Cheng--Duncan--Harvey's Table~2, the following tables give the complete permutation cover $X_0(N)\to X(1)$, its cusps, the action of the subgroup $S$ of exact divisors, and the signature of $\Gamma_0(N)+S$. The signature convention is $(0;e_2,e_3;c)$: genus zero, $e_i$ elliptic points of order $i$, and $c$ cusps. The genus $g_0$ in the first table is that of $X_0(N)$ before taking the $S$-quotient. A cycle expression $a^b$ means $b$ cycles of length $a$.

These are signatures of the effective modular orbifold. The homomorphism $\mathrm{SL}_2(\mathbb Z)\to\mathrm{PSL}_2(\mathbb Z)$ has kernel $\{\pm I\}$, which acts trivially on every dart. Pulling the effective quotient back through this central extension restores the moduli stack of elliptic curves: its generic inertia is $\mu_2$ and its elliptic inertia orders are $4$ and $6$. Thus the displayed finite cover and its unrigidified modular stack are related by that specified central gerbe.
\end{result}

{\small\begin{longtable}{r|r|r|l|l|l|l}
$\ell$ & $\mu$ & $g_0$ & $\sigma_2$ & $\sigma_3$ & cusp widths & $(e_2,e_3;c)$\\\hline
$2$ & 3 & 0 & $1 2$ & $3$ & $1 2$ & $(1,0;2)$\\
$3$ & 4 & 0 & $2^{2}$ & $1 3$ & $1 3$ & $(0,1;2)$\\
$4$ & 6 & 0 & $2^{3}$ & $3^{2}$ & $1^{2} 4$ & $(0,0;3)$\\
$5$ & 6 & 0 & $1^{2} 2^{2}$ & $3^{2}$ & $1 5$ & $(2,0;2)$\\
$6$ & 12 & 0 & $2^{6}$ & $3^{4}$ & $1 2 3 6$ & $(0,0;4)$\\
$7$ & 8 & 0 & $2^{4}$ & $1^{2} 3^{2}$ & $1 7$ & $(0,2;2)$\\
$8$ & 12 & 0 & $2^{6}$ & $3^{4}$ & $1^{2} 2 8$ & $(0,0;4)$\\
$9$ & 12 & 0 & $2^{6}$ & $3^{4}$ & $1^{3} 9$ & $(0,0;4)$\\
$10$ & 18 & 0 & $1^{2} 2^{8}$ & $3^{6}$ & $1 2 5 10$ & $(2,0;4)$\\
$12$ & 24 & 0 & $2^{12}$ & $3^{8}$ & $1^{2} 3^{2} 4 12$ & $(0,0;6)$\\
$13$ & 14 & 0 & $1^{2} 2^{6}$ & $1^{2} 3^{4}$ & $1 13$ & $(2,2;2)$\\
$16$ & 24 & 0 & $2^{12}$ & $3^{8}$ & $1^{4} 4 16$ & $(0,0;6)$\\
$18$ & 36 & 0 & $2^{18}$ & $3^{12}$ & $1^{3} 2^{3} 9 18$ & $(0,0;8)$\\
$25$ & 30 & 0 & $1^{2} 2^{14}$ & $3^{10}$ & $1^{5} 25$ & $(2,0;6)$\\
$6+3$ & 12 & 0 & $2^{6}$ & $3^{4}$ & $1 2 3 6$ & $(2,0;2)$\\
$10+5$ & 18 & 0 & $1^{2} 2^{8}$ & $3^{6}$ & $1 2 5 10$ & $(3,0;2)$\\
$14+7$ & 24 & 1 & $2^{12}$ & $3^{8}$ & $1 2 7 14$ & $(4,0;2)$\\
$18+9$ & 36 & 0 & $2^{18}$ & $3^{12}$ & $1^{3} 2^{3} 9 18$ & $(2,0;4)$\\
$22+11$ & 36 & 2 & $2^{18}$ & $3^{12}$ & $1 2 11 22$ & $(6,0;2)$\\
$30+15$ & 72 & 3 & $2^{36}$ & $3^{24}$ & $1 2 3 5 6 10 15 30$ & $(8,0;4)$\\
$46+23$ & 72 & 5 & $2^{36}$ & $3^{24}$ & $1 2 23 46$ & $(12,0;2)$\\
$12+4$ & 24 & 0 & $2^{12}$ & $3^{8}$ & $1^{2} 3^{2} 4 12$ & $(0,0;4)$\\
$30+6,10,15$ & 72 & 3 & $2^{36}$ & $3^{24}$ & $1 2 3 5 6 10 15 30$ & $(6,0;2)$\\
\end{longtable}
}

\begin{proof}
Use primitive row vectors $(x,y)$ modulo $N$, modulo multiplication by a unit. These are the cosets $\Gamma_0(N)\backslash\mathrm{SL}_2(\mathbb Z)$, by sending a matrix to its bottom row. Its stabilizer is exactly $\Gamma_0(N)$; lifting a primitive residue row to a primitive integral row and completing it to a determinant-one matrix proves surjectivity. The matrices
\[
 s=\begin{pmatrix}0&-1\\1&0\end{pmatrix},\quad
 r=\begin{pmatrix}0&-1\\1&1\end{pmatrix},\quad
 t=\begin{pmatrix}1&1\\0&1\end{pmatrix}
\]
act on the right, with $s^2=r^3=1$ and $sr=t$ in the projective action. Their cycles give the passport. Their numbers of fixed darts give $e_2,e_3$; the $t$-cycles give the cusp widths. Riemann--Hurwitz gives
\[
 g_0=1+\mu/12-e_2/4-e_3/3-c/2.
\]
The executable \texttt{d\_lambencies.py} constructs every primitive row, canonicalizes its unit orbit, and stores all permutations, rather than only their cycle partitions.

Write $f_d=\gcd(d,N/d)$. The cusp label $c_{d,a}$ denotes the class of $a'/d$, where $d\mid N$, $a\in(\mathbb Z/f_d)^\times$, and $a'\equiv a\pmod{f_d}$ is chosen coprime to $d$. For $f_d=1$ we write $a=1$. In particular $c_{1,1}=0$ and $c_{N,1}=\infty$. A reduced rational cusp $x/y$ has label
\[
 d=\gcd(y,N),\qquad a=x(y/d)\pmod{f_d}.
\]
For $\left(\begin{smallmatrix}A&B\\NC&D\end{smallmatrix}\right)\in\Gamma_0(N)$, its transformed numerator is $Ax$ modulo $f_d$, its denominator divided by $d$ is $D(y/d)$ modulo $f_d$, and $AD\equiv1$. This proves invariance. Conversely, completing the two primitive pairs to unimodular matrices and solving the lower-left congruence by B\'ezout gives equivalence exactly for equal $d,a$; equivalently this is the usual cusp criterion used in the opened modular-curve source \cite{RouseWebb2015}. Its width is obtained from the lower-left entry of $\gamma t^w\gamma^{-1}$: the least positive $w$ with $N\mid d^2w$ is
\[
 w_d=N/\gcd(d^2,N).
\]
The Galois action in \cite[Theorem~16]{RouseWebb2015} sends the numerator invariant to its product with a cyclotomic unit (or its inverse, depending on the chosen action convention). Since that action on $(\mathbb Z/f_d)^\times$ is free and transitive, the exact residue field is $\mathbb Q(\zeta_{f_d})$. This also supplies a direct derivation of the field, rather than merely a field over which the cusp becomes rational.

For an exact divisor $e$, choose integers $b,d_0$ with $ed_0-(N/e)b=1$. The matrix
\[
 W_e=\begin{pmatrix}e&b\\N&ed_0\end{pmatrix},\qquad\det W_e=e,
\]
acts on rational cusps; scaling it by $1/\sqrt e$ gives the standard Atkin--Lehner coset. Reduce the resulting fraction and apply the preceding cusp formula. The program checks all products $W_eW_f=W_{ef/\gcd(e,f)^2}$ on cusps. The denominator changes by $d\mapsto ed/\gcd(e,d)^2$, while the table retains the numerator action as well. The latter matters at $N=18$: $W_9$ exchanges the two cusps at denominator $3$ and the two at denominator $6$. The exact cusp fields on a quotient are fixed fields of the stabilizers of these finite Galois orbit sets. In every nontrivial $S$-quotient in this table all cusps are rational. This includes the two merged conjugate pairs for $18+9$.

D.15 below proves genus zero directly from the eta divisor and proves $S$-invariance. For the added cosets, the normalized trace is $\sqrt e(a+d)$. For $e\geq4$ an elliptic element must have trace zero, hence order two. In the sole $e=3$ case, $N=6$, the determinant equation $3ad-2bc=1$ makes $a,d$ odd, so again a trace of absolute value less than two must be zero. The identity coset has no order-three elliptic points for any of the nine nontrivial $S$-rows, as the explicit dart permutations verify. Thus only order two occurs. Hyperbolic area divided by $|S|$ gives
\[
 \frac{\mu}{6|S|}=-2+c+e_2/2,
\]
which proves every remaining signature entry.
\end{proof}

\begin{result}{D.15}{Proved}{The exact tripod comparison and the eta-product check}
The normalization of $X_0(N)\times_{X(1)}X(2)$ has one component of degree $\mu$ over $X(2)$ for odd $N$, and three components, each of degree $\mu/3$, for even $N$. They are defined over $\mathbb Q$. The component containing the standard cusp is isomorphic to $X_0(4N)$ for odd $N$ and to $X_0(2N)$ for even $N$; the other components differ by permutation of the three tripod punctures. For $4\mid N$ there is also the direct cover $X_0(N)\to X_0(4)$ of degree $\mu/6$.

For all twenty-three rows the exact modular function $T^X=1/\eta_{\pi^X}$ is $S$-invariant and descends to a $\mathbb Q$-coordinate on $X_0(N)/S$, with a unique simple pole and expansion $q^{-1}-d^X+O(q)$. The attached computation supplies the full integral cusp divisor and coefficients through $q^{30}$. The assertions of genus zero and invariance follow from those exact divisors and modular transformation criteria, without an inference from a truncated series.
\end{result}
\begin{proof}
The subgroup $\Gamma(2)$ is generated projectively by $t^2$ and $u_2=\left(\begin{smallmatrix}1&0\\-2&1\end{smallmatrix}\right)$. Its orbits on the dart set are the components of the normalized base change, and the three cusp permutations are $t^2,u_2,(t^2u_2)^{-1}$. For odd $N$, the Chinese remainder theorem makes the image of $\Gamma(2)$ act transitively on primitive rows modulo $N$. For even $N$, reduction modulo $2$ records three primitive row classes, and the same lifting argument makes each class an orbit. All three choices are defined over $\mathbb Q$, since the full two-torsion labeling defining $X(2)$ is defined over $\mathbb Q$. Conjugation by $D=\operatorname{diag}(2,1)$ gives, by direct matrix multiplication,
\[
 D\Gamma_0(4N)D^{-1}=\Gamma(2)\cap\Gamma_0(N)\quad(N\text{ odd}),
\]
\[
 D\Gamma_0(2N)D^{-1}=\Gamma(2)\cap\Gamma_0(N)\quad(N\text{ even}).
\]
This conjugation is the algebraic two-isogeny map with the induced two-torsion labeling, so defines the asserted comparison over $\mathbb Q$. The other components are obtained from the rational $S_3$-action on $X(2)$. For the direct cover use $\Gamma_0(N)\subset\Gamma_0(4)$ and the generators $t,u_4=\left(\begin{smallmatrix}1&0\\-4&1\end{smallmatrix}\right)$. Restricting their dart permutations to the orbit of $(0,1)$ computes that cover and its three-point passport. The direct passports for $N=4,8,12,16$ are respectively
\[
 (1,1,1),\quad(1^2,2,2),\quad(1\,3,1\,3,1\,3),\quad(1^4,4,4).
\]
All component permutations, not just their lengths, are in the JSON certificate.

Let $r_\delta$ be the eta exponents of $T^X$, so $r_\delta=-k_\delta$ for the displayed Frame shape. The opened Newman criterion and Ligozat formula in \cite[p.~2]{RouseWebb2015} apply: the program checks
\[
 \sum r_\delta=0,\quad\sum\delta r_\delta\equiv
 \sum(N/\delta)r_\delta\equiv0\pmod{24},\quad
 \prod\delta^{r_\delta}\in(\mathbb Q^\times)^2.
\]
Thus $T^X$ is a weight-zero function with trivial character on $\Gamma_0(N)$ and no zeros or poles in the upper half-plane. At $c_{d,a}$ its order is the exact rational sum
\[
 \operatorname{ord}_{c_{d,a}}T^X=
 \frac N{24}\sum_{\delta\mid N}
 \frac{\gcd(d,\delta)^2r_\delta}{\gcd(d,N/d)d\delta}.
\]
Every value is integral. For every row the negative part consists exactly of the $|S|$ cusps in the orbit of infinity, each with order $-1$; the whole divisor is invariant under the computed $S$-permutations. Consequently $(T^X\circ W_e)/T^X$ is a nonzero constant.

Here are exact constant-one certificates. For each D-type $N=2e$, the matrix
\[
 \begin{pmatrix}e&-(e+1)/2\\2e&-e\end{pmatrix}
\]
has determinant $e$ and trace zero. Its fixed point lies in the upper half-plane, where $T^X$ is finite and nonzero, forcing the constant to be one. At $(N,e)=(30,6)$ choose the trace-zero matrix $(12,-5;30,-12)$, and at $(30,15)$ choose $(15,-8;30,-15)$; their determinant checks are $6$ and $15$. These generate the group containing $W_{10}=W_6W_{15}$. For $(N,e)=(12,4)$, the cusp $c_{2,1}$ is fixed and has eta order zero, again forcing the constant to be one. These certificates are also checked by the program. Therefore $T^X$ descends to the quotient and has one simple pole there. A nonconstant meromorphic function of degree one identifies that smooth projective quotient with $\mathbb P^1$.

The $q$-product has leading order $-1$ and leading coefficient one; its next coefficient is $k_1=-d^X$, as checked against every root system of the opened CDH Table~2 \cite{CDH2014}. Rational $q$-expansions on the canonical rational modular curve and rationality of the Atkin--Lehner maps give the coordinate over $\mathbb Q$. This independently verifies the content of CDH Proposition~2.1 for the table used here.

The quotient maps retain a precise limitation on these covers. For every nontrivial $e$, the function $j\circ W_e$ has pole order $e$ at infinity, while $j$ has pole order one. Hence $j$ does not factor through $X_0(N)/\langle W_e\rangle$. In the special direct-cover row $12+4$, the map $X_0(12)\to X_0(4)$ sends infinity to infinity but sends $W_4(\infty)=c_{3,1}$ to the cusp zero. It therefore does not factor through that quotient either. The explicit covers above, together with the stated quotient maps, are the resulting tripod form; no descent of those particular maps is asserted.
\end{proof}

The next tables make each cusp field and every nontrivial $S$-action explicit. A set in the second subscript lists the distinct numerator invariants; it is not an unreduced rational fraction. The fields with $f=1,2$ are $\mathbb Q$. In the tripod table the other even-level components have the displayed passport with the three punctures permuted. The JSON and the full readable certificate \texttt{code/d\_lambencies.md} additionally list every quotient orbit and every individual eta order.

{\small\begin{longtable}{r|p{0.79\textwidth}}
$N$ & Cusps $c_{d,a}$, with $(w;f)$ meaning width $w$ and field $\mathbb Q(\zeta_f)$\\\hline
2 & $c_{1,\{1\}}:(2;1)$;\ $c_{2,\{1\}}:(1;1)$\\
3 & $c_{1,\{1\}}:(3;1)$;\ $c_{3,\{1\}}:(1;1)$\\
4 & $c_{1,\{1\}}:(4;1)$;\ $c_{2,\{1\}}:(1;2)$;\ $c_{4,\{1\}}:(1;1)$\\
5 & $c_{1,\{1\}}:(5;1)$;\ $c_{5,\{1\}}:(1;1)$\\
6 & $c_{1,\{1\}}:(6;1)$;\ $c_{2,\{1\}}:(3;1)$;\ $c_{3,\{1\}}:(2;1)$;\ $c_{6,\{1\}}:(1;1)$\\
7 & $c_{1,\{1\}}:(7;1)$;\ $c_{7,\{1\}}:(1;1)$\\
8 & $c_{1,\{1\}}:(8;1)$;\ $c_{2,\{1\}}:(2;2)$;\ $c_{4,\{1\}}:(1;2)$;\ $c_{8,\{1\}}:(1;1)$\\
9 & $c_{1,\{1\}}:(9;1)$;\ $c_{3,\{1,2\}}:(1;3)$;\ $c_{9,\{1\}}:(1;1)$\\
10 & $c_{1,\{1\}}:(10;1)$;\ $c_{2,\{1\}}:(5;1)$;\ $c_{5,\{1\}}:(2;1)$;\ $c_{10,\{1\}}:(1;1)$\\
12 & $c_{1,\{1\}}:(12;1)$;\ $c_{2,\{1\}}:(3;2)$;\ $c_{3,\{1\}}:(4;1)$;\ $c_{4,\{1\}}:(3;1)$;\ $c_{6,\{1\}}:(1;2)$;\ $c_{12,\{1\}}:(1;1)$\\
13 & $c_{1,\{1\}}:(13;1)$;\ $c_{13,\{1\}}:(1;1)$\\
14 & $c_{1,\{1\}}:(14;1)$;\ $c_{2,\{1\}}:(7;1)$;\ $c_{7,\{1\}}:(2;1)$;\ $c_{14,\{1\}}:(1;1)$\\
16 & $c_{1,\{1\}}:(16;1)$;\ $c_{2,\{1\}}:(4;2)$;\ $c_{4,\{1,3\}}:(1;4)$;\ $c_{8,\{1\}}:(1;2)$;\ $c_{16,\{1\}}:(1;1)$\\
18 & $c_{1,\{1\}}:(18;1)$;\ $c_{2,\{1\}}:(9;1)$;\ $c_{3,\{1,2\}}:(2;3)$;\ $c_{6,\{1,2\}}:(1;3)$;\ $c_{9,\{1\}}:(2;1)$;\ $c_{18,\{1\}}:(1;1)$\\
22 & $c_{1,\{1\}}:(22;1)$;\ $c_{2,\{1\}}:(11;1)$;\ $c_{11,\{1\}}:(2;1)$;\ $c_{22,\{1\}}:(1;1)$\\
25 & $c_{1,\{1\}}:(25;1)$;\ $c_{5,\{1,2,3,4\}}:(1;5)$;\ $c_{25,\{1\}}:(1;1)$\\
30 & $c_{1,\{1\}}:(30;1)$;\ $c_{2,\{1\}}:(15;1)$;\ $c_{3,\{1\}}:(10;1)$;\ $c_{5,\{1\}}:(6;1)$;\ $c_{6,\{1\}}:(5;1)$;\ $c_{10,\{1\}}:(3;1)$;\ $c_{15,\{1\}}:(2;1)$;\ $c_{30,\{1\}}:(1;1)$\\
46 & $c_{1,\{1\}}:(46;1)$;\ $c_{2,\{1\}}:(23;1)$;\ $c_{23,\{1\}}:(2;1)$;\ $c_{46,\{1\}}:(1;1)$\\
\end{longtable}
\begin{longtable}{r|p{0.78\textwidth}}
$\ell$ & Atkin--Lehner permutations; omitted cusp labels are fixed\\\hline
$6+3$ & $W_{3}=(c_{1,1}\ c_{3,1})(c_{2,1}\ c_{6,1})$\\
$10+5$ & $W_{5}=(c_{1,1}\ c_{5,1})(c_{2,1}\ c_{10,1})$\\
$14+7$ & $W_{7}=(c_{1,1}\ c_{7,1})(c_{2,1}\ c_{14,1})$\\
$18+9$ & $W_{9}=(c_{1,1}\ c_{9,1})(c_{2,1}\ c_{18,1})(c_{3,1}\ c_{3,2})(c_{6,1}\ c_{6,2})$\\
$22+11$ & $W_{11}=(c_{1,1}\ c_{11,1})(c_{2,1}\ c_{22,1})$\\
$30+15$ & $W_{15}=(c_{1,1}\ c_{15,1})(c_{2,1}\ c_{30,1})(c_{3,1}\ c_{5,1})(c_{6,1}\ c_{10,1})$\\
$46+23$ & $W_{23}=(c_{1,1}\ c_{23,1})(c_{2,1}\ c_{46,1})$\\
$12+4$ & $W_{4}=(c_{1,1}\ c_{4,1})(c_{3,1}\ c_{12,1})$\\
$30+6,10,15$ & $W_{6}=(c_{1,1}\ c_{6,1})(c_{2,1}\ c_{3,1})(c_{5,1}\ c_{30,1})(c_{10,1}\ c_{15,1})$\newline $W_{10}=(c_{1,1}\ c_{10,1})(c_{2,1}\ c_{5,1})(c_{3,1}\ c_{30,1})(c_{6,1}\ c_{15,1})$\newline $W_{15}=(c_{1,1}\ c_{15,1})(c_{2,1}\ c_{30,1})(c_{3,1}\ c_{5,1})(c_{6,1}\ c_{10,1})$\\
\end{longtable}
\begin{longtable}{r|r|r|l|l}
$N$ & components & degree & passport of one component over $X(2)$ & direct degree to $X_0(4)$\\\hline
2 & 3 & 1 & $(1,1,1)$ & --\\
3 & 1 & 4 & $(1\,3,1\,3,1\,3)$ & --\\
4 & 3 & 2 & $(1^{2},2,2)$ & 1\\
5 & 1 & 6 & $(1\,5,1\,5,1\,5)$ & --\\
6 & 3 & 4 & $(1\,3,1\,3,1\,3)$ & --\\
7 & 1 & 8 & $(1\,7,1\,7,1\,7)$ & --\\
8 & 3 & 4 & $(1^{4},4,4)$ & 2\\
9 & 1 & 12 & $(1^{3}\,9,1^{3}\,9,1^{3}\,9)$ & --\\
10 & 3 & 6 & $(1\,5,1\,5,1\,5)$ & --\\
12 & 3 & 8 & $(1^{2}\,3^{2},2\,6,2\,6)$ & 4\\
13 & 1 & 14 & $(1\,13,1\,13,1\,13)$ & --\\
14 & 3 & 8 & $(1\,7,1\,7,1\,7)$ & --\\
16 & 3 & 8 & $(1^{4}\,2^{2},8,8)$ & 4\\
18 & 3 & 12 & $(1^{3}\,9,1^{3}\,9,1^{3}\,9)$ & --\\
22 & 3 & 12 & $(1\,11,1\,11,1\,11)$ & --\\
25 & 1 & 30 & $(1^{5}\,25,1^{5}\,25,1^{5}\,25)$ & --\\
30 & 3 & 24 & $(1\,3\,5\,15,1\,3\,5\,15,1\,3\,5\,15)$ & --\\
46 & 3 & 24 & $(1\,23,1\,23,1\,23)$ & --\\
\end{longtable}
\begin{longtable}{r|l|r|r}
$\ell$ & $\pi^X$ (positive numerator, negative denominator) & $[q^0]T^X$ & $[q]T^X$\\\hline
$2$ & $\frac{2^{24}}{1^{24}}$ & -24 & 276\\
$3$ & $\frac{3^{12}}{1^{12}}$ & -12 & 54\\
$4$ & $\frac{4^{8}}{1^{8}}$ & -8 & 20\\
$5$ & $\frac{5^{6}}{1^{6}}$ & -6 & 9\\
$6$ & $\frac{2\,6^{5}}{1^{5}\,3}$ & -5 & 6\\
$7$ & $\frac{7^{4}}{1^{4}}$ & -4 & 2\\
$8$ & $\frac{2^{2}\,8^{4}}{1^{4}\,4^{2}}$ & -4 & 4\\
$9$ & $\frac{9^{3}}{1^{3}}$ & -3 & 0\\
$10$ & $\frac{2\,10^{3}}{1^{3}\,5}$ & -3 & 1\\
$12$ & $\frac{2^{2}\,3\,12^{3}}{1^{3}\,4\,6^{2}}$ & -3 & 2\\
$13$ & $\frac{13^{2}}{1^{2}}$ & -2 & -1\\
$16$ & $\frac{2\,16^{2}}{1^{2}\,8}$ & -2 & 0\\
$18$ & $\frac{2\,3\,18^{2}}{1^{2}\,6\,9}$ & -2 & 0\\
$25$ & $\frac{25}{1}$ & -1 & -1\\
$6+3$ & $\frac{2^{6}\,6^{6}}{1^{6}\,3^{6}}$ & -6 & 15\\
$10+5$ & $\frac{2^{4}\,10^{4}}{1^{4}\,5^{4}}$ & -4 & 6\\
$14+7$ & $\frac{2^{3}\,14^{3}}{1^{3}\,7^{3}}$ & -3 & 3\\
$18+9$ & $\frac{2^{3}\,3^{2}\,18^{3}}{1^{3}\,6^{2}\,9^{3}}$ & -3 & 3\\
$22+11$ & $\frac{2^{2}\,22^{2}}{1^{2}\,11^{2}}$ & -2 & 1\\
$30+15$ & $\frac{2^{2}\,3\,5\,30^{2}}{1^{2}\,6\,10\,15^{2}}$ & -2 & 1\\
$46+23$ & $\frac{2\,46}{1\,23}$ & -1 & 0\\
$12+4$ & $\frac{2^{4}\,3^{4}\,12^{4}}{1^{4}\,4^{4}\,6^{4}}$ & -4 & 6\\
$30+6,10,15$ & $\frac{2^{3}\,3^{3}\,5^{3}\,30^{3}}{1^{3}\,6^{3}\,10^{3}\,15^{3}}$ & -3 & 3\\
\end{longtable}
}


\subsection{The reconstruction, its stage, and the five link diagrams}

\begin{result}{D.16}{Proved}{Canonical reconstruction from the supplied face and collar diagram}
Let $\mathcal V$ be the real oriented blow-up of a smooth complex normal-crossings Deligne--Mumford pair $(\mathcal X,\mathcal D)$, with its full deployed face diagram, collars, and maps to the strata. Collapsing precisely the normal angular circles over the zero radial coordinates recovers $\mathcal X$ as a topological stack. Equivalently, the mapping-cylinder pushout
\[
 \widehat X=X^*\mathop{\amalg}_{\dot X} C(p)
\]
of cote~150, scans~32--35, performed successively on the deployed strata, gives this same stack. For the supplied data, the final isomorphism is canonical and independent of the order of a compatible induction.

For $\overline{\mathcal M}_{g,n}$, the fundamental groupoid is already that of the compactification after all open divisors have been attached to the open stratum. In the deepest-strata-first reconstruction starting from $\mathcal V$, it is already that groupoid after the first layer, consisting of all maximal-corner fibres. In both descriptions it is the connected simply connected groupoid. These are two explicitly related filtrations of the same final quotient.
\end{result}
\begin{proof}
On a uniformizing chart, write the boundary equation as $z_1\cdots z_q=0$. The blow-up and its blow-down are
\[
 (\rho_i,\theta_i,z_{q+1},\ldots,z_m)\longmapsto
 (\rho_i e^{\mathrm i\theta_i},z_{q+1},\ldots,z_m),\qquad \rho_i\geq0.
\]
Its fibres identify exactly the $\theta_i$ for which $\rho_i=0$. The resulting quotient chart is the original complex chart: polar coordinates give a continuous bijection from the compact local quotient to its Hausdorff image, and the same assertion on relatively compact subcharts gives a homeomorphism. The construction is equivariant for the finite chart group and for coordinate changes. Taking the quotient groupoids therefore gives the original topological stack.

For a closed attaching stratum $S$, let $p:L_S\to S$ be its deployed normal link. Its mapping cylinder is
\[
 C(p)=\bigl(L_S\times[0,1]\amalg S\bigr)/((\ell,0)\sim p(\ell)).
\]
The collar identifies $L_S\times\{1\}$ with the attaching boundary of the previously cut piece. Gluing this cylinder implements exactly the preceding angular quotient on the corresponding normal cone. Repeating on all faces recovers the common quotient of $\mathcal V$. The coordinate equivalence relations commute: collapsing coordinate $i$ and then $j$ identifies exactly the same pairs as collapsing $j$ and then $i$. More generally the universal property of the quotient by the union of these relations identifies every compatible iterated pushout with one fixed quotient. The induced isomorphisms are unique over that quotient and satisfy their cocycle law. This is the asserted canonical independence for the given face/collar data; it requires no claim that collars have a preferred choice.

For the first filtration, set $\mathcal X^{\leq1}=\mathcal M_{g,n}\cup\{\text{open boundary divisors}\}$. Its inclusion into $\overline{\mathcal M}_{g,n}$ induces an isomorphism on $\pi_1$: the omitted strata have complex codimension at least two, and a loop or a homotopy of loops can be moved off them by real transversality in the uniformizing charts. Van Kampen for attaching an open divisor kills its normal meridian. Every essential simple closed curve gives a boundary degeneration, so the resulting quotient is $\Gamma_{g,n}$ modulo the normal subgroup generated by all Dehn twists. These twists generate the pure mapping class group. In positive genus this is exactly Labru\`ere--Paris, Corollary~2.11(i), with its puncture-induction proof on pages~86--89 \cite{D_LabruereParis2001}; in genus zero it follows from the pure-braid generators given in Part~B. Thus this quotient is trivial.

For the second filtration, the inclusion of the interior into $\mathcal V$ is a homotopy equivalence by pushing inward along a collar, so $\pi_1(\mathcal V)=\Gamma_{g,n}$. Over a maximal boundary stratum the deployed face is a torus bundle over the finite automorphism stack of a pants graph. The map to that stack is surjective on $\pi_1$, with kernel the image of its normal torus lattice. Replacing it by its mapping cylinder therefore imposes only the relations killing that lattice. Every essential curve extends to a pants decomposition; the normal closure of the lattices over all maximal strata is consequently the subgroup generated by all Dehn twists. The first complete layer has trivial fundamental group. Subsequent angular collapses again have connected fibres and induce quotients on $\pi_1$, so the group stays trivial. These intermediate quotients need not themselves be smooth stacks; they are the topological pushouts specified by the reconstruction. The two filtrations map to the same final quotient just proved.
\end{proof}

\paragraph{The four reconstruction cases.}
For $(0,4)$, $\mathcal V$ is a pair of pants. The three face maps are $S^1\to\{\text{point}\}$, and their mapping cylinders are disks. The pushout is $\mathbb P^1$. Killing the three peripheral loops already gives its trivial fundamental group.

For $(0,5)$, the ten deployed codimension-one faces are the normal circle bundles over the cut versions of the ten boundary copies of $\mathbb P^1$; their open bases are tripods. Fifteen deployed corner fibres are two-tori, one for each edge of the Petersen incidence graph. Each of the ten normal line bundles has degree $-1$. On a corner the two collapse maps are $(e^{\mathrm i\theta_1},e^{\mathrm i\theta_2})\mapsto e^{\mathrm i\theta_1}$ and $(e^{\mathrm i\theta_1},e^{\mathrm i\theta_2})\mapsto e^{\mathrm i\theta_2}$, followed by the corresponding remaining circle collapse. Their common quotient is the local chart $(\rho_1e^{\mathrm i\theta_1},\rho_2e^{\mathrm i\theta_2})$. Gluing these charts with the prescribed degree-$-1$ normal bundles recovers the degree-five del Pezzo surface $\overline M_{0,5}$. Collapsing all fifteen maximal two-torus fibres is already the first simply connected stage.

For $(1,1)$, the cut space is the compactified modular orbifold with its cusp cut open. The cusp face is $S^1\times B\mu_2\to B\mu_2$: the involution interchanges the two branches of the nodal rational curve and acts trivially on their tensor-product smoothing parameter. Its mapping cylinder fills the cusp while retaining that inertia. The resulting analytic stack is $\overline{\mathcal M}_{1,1}\simeq\mathbb P(4,6)$, including the interior inertia of orders $4$ and $6$ and generic inertia $\mu_2$. The boundary meridian maps to the usual Dehn twist $T\in\mathrm{SL}_2(\mathbb Z)$. The presentation $S^4=1$, $S^2=(ST)^3$ shows directly that setting $T=1$ kills the whole group. Thus the cusp attachment gives the compactification's trivial fundamental group; it retains the stack inertia.

For $(1,2)$, the two divisor normalizations are
\[
 [\overline M_{0,4}/\mathfrak S_2]\quad\hbox{and}\quad
 \overline{\mathcal M}_{1,1}\times\overline M_{0,3}.
\]
In the first, $\mathfrak S_2$ exchanges the node branches and fixes the two marked legs. The two maximal graph types are the loop with a rational tail and the two-edge cycle. Their deployed corner fibres are two-torus bundles over their graph-automorphism stacks; their local groups are respectively the groups $D$ and $C$ of B.10. The loop-tail corner has one retained normal direction for each divisor. The cycle corner has two distinct retained-branch incidences with the irreducible divisor, exchanged by its edge-swap automorphism. The face maps forget the smoothed angular coordinates and then use the sewing maps of B.11. Hence both incidences, and their finite automorphisms, occur in the pushout. The local two-coordinate quotient recovers the node-smoothing chart and the global quotient is $\overline{\mathcal M}_{1,2}$. Killing the two maximal meridian lattices kills every Dehn twist by the extension-to-pants argument; this is again the first simply connected stage.

\begin{result}{D.17}{Proved}{The five depth-three diagrams and the role of the four elementary types}
For $(0,5),(0,6),(1,2),(1,3),(2,0)$, the table below gives every boundary stratum of depth at most three, with its cut-component types and graph automorphisms. On the incidence category of deployed flags, the local group and each arrow are given by the following explicit normal-bundle and sewing construction. Its fundamental group is $\pi_1(L_{g,n})$. Consequently it equals $\Gamma_{g,n}$ in the three complexity-three types $(0,6),(1,3),(2,0)$; in the two complexity-two types its map to $\Gamma_{g,n}$ has kernel $F_{\aleph_0}$.

For every stable type, the action complex of groups on $P_2/\Gamma_{g,n}$ has fundamental group $\Gamma_{g,n}$. Its data use strata of dimension at most two. The connected one-dimensional stage supplies generators, the two-dimensional stage supplies relations, and a stratum of dimension three supplies relations among those relations. The four positive-dimensional single-component elementary types are exactly $(0,4),(1,1),(0,5),(1,2)$.
\end{result}
\begin{proof}
Here is the full geometric specification of the local groups, including their extension classes. For a graph $G$, let
\[
 \mathcal S_G=\left[\prod_v\mathcal M_{g_v,n_v}/\operatorname{Aut}(G)\right].
\]
For each edge $e=\{h,h'\}$ its normal line is
$\mathcal N_e=L_h^\vee\otimes L_{h'}^\vee$, of first Chern class
$-\psi_h-\psi_{h'}$. Let $\mathcal T_G$ be the simultaneous unit-circle bundle of these lines, with the graph group permuting the factors. Define $H_G=\pi_1(\mathcal T_G)$. Pulling to the product of cut-surface Teichm\"uller spaces makes this bundle trivial and proves
\[
 1\longrightarrow\mathbb Z^{E(G)}
 \longrightarrow H_G
 \longrightarrow\pi_1(\mathcal S_G)\longrightarrow1.
\]
The bundle itself specifies the extension, including its graph-automorphism lifts. Equivalently $H_G$ is the stabilizer of the chosen multicurve under sewing; opposite boundary twists of a cut edge are identified with the same global node twist.

A flag colours the edges by the stage at which they are contracted. Its group is the subgroup of $H_G$ preserving that colouring. The two edge injections into adjacent vertex groups are induced by sewing the specified boundary pairs and by retaining the other parametrized normal circles. On the node lattices they send a retained edge basis vector to that same edge basis vector. Graph automorphisms are restricted to those preserving the flag, before sewing. Thus multiple incidences are retained even when their contracted graph types coincide. All graph lists, individual contractions, and edge multiplicities are in the supplied JSON computation; the displayed table groups only isomorphic label-permutation patterns.

For completeness, the twisting elements in an ordinary presentation are specified as follows. Take the barycentric face category of the action complex, choose a representative of each orbit of a chain, and for an arrow $\alpha:A\to B$ choose $g_\alpha\in\Gamma_{g,n}$ carrying the chosen representative of $B$ to its designated face of the chosen representative of $A$. Then
\[
 \psi_\alpha(x)=g_\alpha^{-1}xg_\alpha,\qquad
 f_{\beta,\alpha}=g_\beta^{-1}g_\alpha^{-1}g_{\beta\alpha}
 \in H_C\quad(A\xrightarrow{\alpha}B\xrightarrow{\beta}C).
\]
The equality $\psi_\beta\psi_\alpha=\operatorname{Ad}(f_{\beta,\alpha})\psi_{\beta\alpha}$ follows by multiplying these three words. Multiplying four transports gives the cocycle identity, so this specifies the twisting on every triangle. The presentation has all local groups and arrow letters $u_\alpha$, the relations
\[
 u_\alpha x u_\alpha^{-1}=\psi_\alpha(x),\qquad
 u_\beta u_\alpha=f_{\beta,\alpha}u_{\beta\alpha},
\]
and $u_\alpha=1$ on a chosen maximal tree. This construction is the Borel action groupoid of the actual chain complex: equivariant simplexwise contraction gives its classifying-space comparison, so its homotopy fibre over $B\Gamma_{g,n}$ is that chain complex. For the depth-three link diagram D.2 supplies the one-type equivalence with $L_{g,n}$; D.1 proves exactly the group conclusions stated above.

For $P_2$, D.5 gives simple connectedness, and the same Borel fibre sequence gives $\pi_1=\Gamma_{g,n}$. If $d\leq2$, this assertion includes the empty multicurve, whose vertex is the open stratum. The exact map from the nonempty boundary diagram is induced by adjoining that empty multicurve and all incident chains; its kernel in the two complexity-two examples is the free group in D.1. Thus the low-dimensional and punctured-boundary presentations are related by this concrete map.

At a stratum of dimension one, the positive-complexity cut components have type $(0,4)$ or $(1,1)$, with pants factors. At dimension two they are either a single $(0,5)$ or $(1,2)$ component, or two complexity-one components. This follows by summing $3g_v-3+n_v$ and imposing stability, and proves that precisely the four elementary types furnish the local generator and relation data. D.5's connectedness of $P_1$ and simple connectedness of $P_2$ prove their global roles. At residual dimension three, the upper link is a join of curve complexes with total residual complexity three; its homotopy type is a wedge of two-spheres. Coning that link attaches three-dimensional cells up to homotopy and specifies identities among the two-dimensional relations. The single positive-complexity types are $(0,6),(1,3)$ and, only for the entire closed surface, $(2,0)$; the other possibilities are one complexity-two and one complexity-one component, or three complexity-one components. This lists every stratum type. For the three requested dimension-three moduli spaces, the sole dimension-three stratum is the open stratum itself. The two requested dimension-two spaces have no dimension-three stratum.
\end{proof}

The following enumeration was executed. Vertices are numbered from zero; $00$ denotes a loop and repeated entries denote parallel edges. The labelled number counts the individual strata in that label-permutation orbit. The normal-bundle rule, with these exact edges and cut types, specifies each vertex group and every coloured-flag subgroup.

{\small\paragraph{Type $(0,5)$.}
\begin{longtable}{c|r|l|l|l|r}
$q$ & labelled number & genera; leg counts & edge multiset & cut types & $|\operatorname{Aut}G|$\\\hline
1 & 10 & $(0,0);(2,3)$ & $\{01\}$ & $(0,3),(0,4)$ & 1\\
2 & 15 & $(0,0,0);(1,2,2)$ & $\{01,02\}$ & $(0,3),(0,3),(0,3)$ & 1\\
\end{longtable}
\paragraph{Type $(0,6)$.}
\begin{longtable}{c|r|l|l|l|r}
$q$ & labelled number & genera; leg counts & edge multiset & cut types & $|\operatorname{Aut}G|$\\\hline
1 & 15 & $(0,0);(2,4)$ & $\{01\}$ & $(0,3),(0,5)$ & 1\\
1 & 10 & $(0,0);(3,3)$ & $\{01\}$ & $(0,4),(0,4)$ & 1\\
2 & 60 & $(0,0,0);(1,2,3)$ & $\{01,02\}$ & $(0,3),(0,3),(0,4)$ & 1\\
2 & 45 & $(0,0,0);(2,2,2)$ & $\{01,02\}$ & $(0,3),(0,3),(0,4)$ & 1\\
3 & 15 & $(0,0,0,0);(0,2,2,2)$ & $\{01,02,03\}$ & $(0,3),(0,3),(0,3),(0,3)$ & 1\\
3 & 90 & $(0,0,0,0);(1,1,2,2)$ & $\{01,02,13\}$ & $(0,3),(0,3),(0,3),(0,3)$ & 1\\
\end{longtable}
\paragraph{Type $(1,2)$.}
\begin{longtable}{c|r|l|l|l|r}
$q$ & labelled number & genera; leg counts & edge multiset & cut types & $|\operatorname{Aut}G|$\\\hline
1 & 1 & $(0);(2)$ & $\{00\}$ & $(0,4)$ & 2\\
1 & 1 & $(0,1);(2,0)$ & $\{01\}$ & $(0,3),(1,1)$ & 1\\
2 & 1 & $(0,0);(0,2)$ & $\{00,01\}$ & $(0,3),(0,3)$ & 2\\
2 & 1 & $(0,0);(1,1)$ & $\{01,01\}$ & $(0,3),(0,3)$ & 2\\
\end{longtable}
\paragraph{Type $(1,3)$.}
\begin{longtable}{c|r|l|l|l|r}
$q$ & labelled number & genera; leg counts & edge multiset & cut types & $|\operatorname{Aut}G|$\\\hline
1 & 1 & $(0);(3)$ & $\{00\}$ & $(0,5)$ & 2\\
1 & 3 & $(0,1);(2,1)$ & $\{01\}$ & $(0,3),(1,2)$ & 1\\
1 & 1 & $(0,1);(3,0)$ & $\{01\}$ & $(0,4),(1,1)$ & 1\\
2 & 1 & $(0,0);(0,3)$ & $\{00,01\}$ & $(0,3),(0,4)$ & 2\\
2 & 3 & $(0,0);(1,2)$ & $\{00,01\}$ & $(0,3),(0,4)$ & 2\\
2 & 3 & $(0,0);(1,2)$ & $\{01,01\}$ & $(0,3),(0,4)$ & 2\\
2 & 3 & $(0,0,1);(1,2,0)$ & $\{01,02\}$ & $(0,3),(0,3),(1,1)$ & 1\\
3 & 3 & $(0,0,0);(0,1,2)$ & $\{00,01,12\}$ & $(0,3),(0,3),(0,3)$ & 2\\
3 & 3 & $(0,0,0);(0,1,2)$ & $\{01,01,02\}$ & $(0,3),(0,3),(0,3)$ & 2\\
3 & 1 & $(0,0,0);(1,1,1)$ & $\{01,02,12\}$ & $(0,3),(0,3),(0,3)$ & 1\\
\end{longtable}
\paragraph{Type $(2,0)$.}
\begin{longtable}{c|r|l|l|l|r}
$q$ & labelled number & genera; leg counts & edge multiset & cut types & $|\operatorname{Aut}G|$\\\hline
1 & 1 & $(1);(0)$ & $\{00\}$ & $(1,2)$ & 2\\
1 & 1 & $(1,1);(0,0)$ & $\{01\}$ & $(1,1),(1,1)$ & 2\\
2 & 1 & $(0);(0)$ & $\{00,00\}$ & $(0,4)$ & 8\\
2 & 1 & $(0,1);(0,0)$ & $\{00,01\}$ & $(0,3),(1,1)$ & 2\\
3 & 1 & $(0,0);(0,0)$ & $\{00,01,11\}$ & $(0,3),(0,3)$ & 8\\
3 & 1 & $(0,0);(0,0)$ & $\{01,01,01\}$ & $(0,3),(0,3)$ & 12\\
\end{longtable}
}

\paragraph{Concrete group forms in the first three examples.}
For $(0,5)$, each Petersen vertex $i$ has group
\[
 H_i=\langle f_i,h_{ij_1},h_{ij_2},h_{ij_3}\mid
 [f_i,h_{ij_a}]=1,\quad h_{ij_1}h_{ij_2}h_{ij_3}=f_i\rangle.
\]
This is $\mathbb Z\times F_2$; the relation records normal degree $-1$ and fixes the lifts of the three base peripheral loops. The edge group is $\mathbb Z^2$. Plumbing injects its basis as $(f_i,h_{ij})$ at $i$ and $(h_{ji},f_j)$ at $j$. Non-tree edges carry their arrow letters in the preceding presentation. Abelianizing gives the matrix $-I+A$ of D.4, and hence rank $11$, confirming that this link group is the specified extension of $\Gamma_{0,5}$.

For $(0,6)$, the fifteen $2|4$ divisor groups are the normal circle extensions of $\Gamma_{0,5}$, and the ten $3|3$ divisor groups those of $F_2\times F_2$, with Euler classes $-\psi_h-\psi_{h'}$. The $105$ two-edge groups are the two-circle extensions over a tripod, hence abstractly $\mathbb Z^2\times F_2$; their chosen peripheral lifts remain the prescribed normal-line lifts. The $105$ three-edge groups are $\mathbb Z^3$. All flag maps are the specified sewing maps, and the triangles are exactly the nested edge-contraction flags of the enumerated trees.

For $(1,2)$, B.10--B.11 give the completely presented node-only groups
\[
\begin{aligned}
 N&=\langle a,c,r\mid [a,c]=[a,r]=1,\ r^2=1\rangle,\\
 T&=\langle a,b\mid aba=bab\rangle,\\
 D&=\langle a,h\mid[a,h]=1\rangle,\\
 C&=\langle a,c,q\mid[a,c]=1,\ q^2=1,\ qaq^{-1}=c\rangle.
\end{aligned}
\]
Here $N,T$ are the irreducible and separating divisor groups; $D,C$ are the loop-tail and two-cycle groups. Their node lattices are $\langle a\rangle$, $\langle(ab)^6\rangle$, $\langle a,h^2\rangle$, and $\langle a,c\rangle$. The injections $D\to T$ and $D\to N$ are
\[
 a\mapsto a,\ h\mapsto(ab)^3;\qquad
 a\mapsto a,\ h\mapsto ac^{-1}r.
\]
For a branch of $C\to N$ the flag subgroup is $\langle a,c\rangle$ and its injection is the identity; the other branch is its $q$-conjugate. This restriction is necessary because $q$ exchanges the retained and smoothed edges. The map to the open presentation
$\Gamma_{1,2}=B_4(a,b,c)/\langle(abc)^4\rangle$
is
$h=(ab)^3$, $q=(abc)^2$, $r=(aba)^{-1}(abcaba)(aba)$.
The word checks and external-leg-framed versions are proved and executed in B.9--B.11.

The remaining generator-by-generator identification of the entire $P_2$ action presentation with the specific five- and six-point finite presentations is separately recorded as open in B.8/B.15. The geometric action presentation above does not claim to complete that word dictionary.


\subsection{The exact graph, boundary, and tropical comparison}

\begin{result}{D.18}{Proved}{The corrected maquette and its comparison fibres}
Let $\mathsf G_{g,n}^{+}$ be the category of stable labelled graphs with at least one edge, with specified contractions and graph isomorphisms as morphisms. Define
\[
 \mathcal B_p=
 \{\,(G,\lambda):|E(G)|=p+1,\ 
       \lambda:E(G)\xrightarrow{\sim}\{0,\ldots,p\}\,\}/\cong .
\]
Relabelling gives the symmetric-group action; an injection of label sets gives a face by contracting precisely the omitted edges. Then $\mathcal B$ is a symmetric semisimplicial set, its realization is the boundary complex, and the map assigning the label weights to the edge lengths is a homeomorphism
\[
 |\mathcal B|\xrightarrow{\sim}\Delta_{g,n}.
\]
Here $\Delta_{g,n}$ is the space of stable tropical curves of total finite-edge length one. The category with the automorphisms retained has the explicit comparison
\[
 B\mathsf G_{g,n}^{+}\simeq
 \operatorname*{hocolim}_{(\mathsf G_{g,n}^{+})^{\rm op}}\Delta_G
 \longrightarrow
 \operatorname*{colim}_{(\mathsf G_{g,n}^{+})^{\rm op}}\Delta_G
 =\Delta_{g,n},\qquad
 \Delta_G=\{(\ell_e)\geq0:\sum\ell_e=1\}.
\]
The actual point fibre of this homotopy-colimit model over a metric graph $(G,\ell)$ has homotopy type $B\operatorname{Aut}(G,\ell)$. The map need not be a quasifibration; this assertion concerns its actual point fibres. Adding the zero-edge graph makes the full graph category contractible, since that graph is terminal.

The analytic boundary end has a compatible comparison to these models. Over a positive-length metric multicurve $(\sigma,\ell)$, the point fibre of its Borel-model map to $\Delta_{g,n}$ is
$B\operatorname{Stab}_{\Gamma_{g,n}}(\sigma,\ell)$.
The map of that fibre to the graph-model fibre is induced by the graph-action homomorphism, whose kernel is the node-framed pure cut-surface group. In particular its homotopy fibre is the classifying space of that kernel.
\end{result}
\begin{proof}
An edge contraction is specified by its contracted edge set and the induced quotient of vertices, half-edges and genera. Consecutive contractions compose by contracting the union of the two sets, transported through the first quotient. Thus injections of finite label sets give contravariant face operations and satisfy their functor identities exactly. Relabelling commutes with these operations. This constructs $\mathcal B$ and proves all its symmetric semisimplicial identities.

A point in its realization has an edge-labelled graph and nonnegative coordinates of sum one, modulo label permutations and the face relations. Contracting exactly the zero-coordinate edges gives a metric stable graph with positive edge lengths. Conversely, choose an edge labelling of a metric graph and use its lengths as simplex coordinates; changing that choice is precisely the symmetric identification. These operations are inverse. The topology on both sides is the final topology from the finitely many graph simplices, so the bijection and its inverse are continuous. In a node-smoothing chart the branches of the boundary are exactly the edges of the dual graph, and simultaneous normalization marks precisely the chosen branches. The boundary-complex face operation is therefore the same edge contraction, including graph monodromy. This gives the asserted identification with the boundary complex. It is also the construction in Chan--Galatius--Payne, Definitions~2.3--2.4 and~3.1 and Theorem~6.1 \cite{D_CGPMarked2022}; we opened the definitions and proof. The unmarked statement and descent proof are Corollaries~5.6--5.7 of \cite{D_CGP2021}.

Every $\Delta_G$ is nonempty and contractible, so the map from its homotopy colimit to the classifying space of the index category is a homotopy equivalence, by the levelwise bar-construction comparison with a constant point diagram. A contraction induces the inclusion of its target simplex as the corresponding zero-coordinate face, and isomorphisms act by their actual coordinate permutations. The ordinary colimit is therefore exactly the quotient just constructed. Projection from the bar construction to that colimit is the displayed map.

Fix a metric graph $x=(G,\ell)$ with all lengths positive. The fibre of that bar-construction map is the nerve of the category of its lifts: objects are graph-simplex representatives mapping to $x$, and arrows are the corresponding face maps and isomorphisms preserving the metric point. Contract all zero-length edges in a lift. This gives a functor to the groupoid of positive-length representatives of $x$, which is equivalent to the one-object groupoid $\operatorname{Aut}(G,\ell)$. The canonical zero-edge contractions give a natural transformation from its inclusion followed by contraction to the identity of the lift category. Its nerve is thus homotopy equivalent to $B\operatorname{Aut}(G,\ell)$. This proves the fibre calculation, including automorphisms of zero-length refinements, which disappear through that natural transformation. For example, at $(1,1)$ the unique nonempty graph is a loop with its flip automorphism: the map is $BC_2\to\{\text{point}\}$. This computes the failure of the original asserted equality of classifying spaces. The zero-edge graph, on the other hand, has a unique contraction from every graph, proving the full category's contractibility.

For the analytic comparison, represent the end as
$E\Gamma_{g,n}\times_{\Gamma_{g,n}}|\mathcal C_{g,n}|$
as in D.1. Send a weighted multicurve to its dual graph with those weights as edge lengths. The classification of marked cut surfaces shows that the orbits are exactly the metric stable graphs, so the orbit space is the same $\Delta_{g,n}$. The preimage of one orbit in the Borel construction is
$E\Gamma_{g,n}\times_{\Gamma_{g,n}}\Gamma_{g,n}/H\simeq BH$,
where $H=\operatorname{Stab}(\sigma,\ell)$.
Cutting the surface and forgetting its marking gives a functor to the metric-graph lift category and hence the compatible comparison above. On this fibre its homomorphism is
\[
 1\longrightarrow K_G\longrightarrow H
 \longrightarrow\operatorname{Aut}(G,\ell)\longrightarrow1,
 \qquad
 1\longrightarrow\mathbb Z^{E(G)}
 \longrightarrow K_G
 \longrightarrow\prod_v\Gamma_{g_v,n_v}\longrightarrow1.
\]
Surjectivity onto graph automorphisms follows by choosing oriented homeomorphisms between the cut components with the prescribed flag labels and sewing them. The kernel consists exactly of mapping classes fixing each component and each half-edge; the cut-surface sewing sequence gives its second displayed extension. Thus $BK_G$ is the homotopy fibre of the stated map $BH\to B\operatorname{Aut}(G,\ell)$. The normal Euler classes are those in D.17 and A.12.
\end{proof}

The terminology in Lead~6 can now be assigned exact maps. Specified MD-graph generizations, with their graph automorphisms, give $\mathsf G_{g,n}$ and its nerve. Taking isomorphism classes of edge-labelled graphs gives $\mathcal B$. Taking isomorphism classes of strictly contracting flags gives the barycentric subdivision of $\mathcal B$: on each graph simplex this is the ordinary subdivision by nested nonempty coordinate faces, and the construction commutes with every graph isomorphism and contraction. Its realization maps homeomorphically to $|\mathcal B|$ by sending a face-chain vertex to the corresponding simplex barycenter. This describes the comparison without a degreewise identification of edge-labelled cells with contraction chains. The levelwise component sets of the original groupoid-valued flag data are addressed in A.11 and A.13; their lost middle-endpoint identifications need not satisfy a Segal composition law.

For the two examples in the added task, this map has concrete consequences. For the tripod, the tropical boundary is three points; the analytic end's three circle fibres are its three meridian groups, while the trivially valued tube adds a positive radial interval at each point. For $\overline M_{0,5}$, graph automorphisms are trivial and the graph model is the Petersen graph. The analytic end maps to it with circle-bundle groups over vertex interiors and $\mathbb Z^2$ over edge interiors; the nonarchimedean tube retracts to it by D.4. The exact maps, the Betti numbers in D.4, and the proper-modification comparisons in D.3 establish all asserted relationships in these examples.


\subsection{The full low-genus completion range and two source corrections}

\begin{result}{D.19}{Known; transfer proved}{Congruence comparison in genus at most two}
For every stable type of genus $0,1,$ or $2$, the natural map
$\widehat\Gamma_{g,n}\to\check\Gamma_{g,n}$
is an isomorphism. Consequently the full finite-level and procongruence completions of every $P_k$ are canonically isomorphic in these genera. In particular, their $P_1$ are connected and their $P_k$, $k\geq2$, have trivial profinite fundamental group. The full unmodified curve complex is connected with trivial profinite fundamental group in this range exactly when $d\geq3$; the two connected complexity-two exceptions have nontrivial profinite fundamental group by D.12 and D.21.
\end{result}
\begin{proof}
The genus-zero congruence theorem is the puncture induction in Boggi's Proposition~2.6, reproducing Asada's argument: start with $\Gamma_{0,3}=1$, use the completed Birman sequence, and use the tangential lift to identify the point-pushing closure with the full profinite surface group \cite{D_Boggi2009}. The opened Theorem~3.5 and Lemmas~3.6--3.7 of the same paper prove faithfulness for the full hyperelliptic group with at least one puncture. In genera one and two the hyperelliptic locus is the entire moduli stack, so this gives every stable genus-one type and every genus-two type with $n\geq1$. The closed genus-two case requires an additional argument; it is not part of that open-surface theorem.

Here is the independent closed-surface transfer, which avoids both problematic claims discussed in D.20. Write
\[
 P=\widehat{\pi_1(S_2)},\qquad
 Q=\widehat{\pi_1(\operatorname{Conf}_2(S_2))},\qquad
 F=\widehat{\pi_1(S_2\setminus\{x\})}.
\]
The full universal-curve and configuration sequences give
\[
 1\longrightarrow Q\longrightarrow\widehat\Gamma_{2,2}
       \longrightarrow\widehat\Gamma_{2,0}\longrightarrow1,
 \qquad
 1\longrightarrow F\longrightarrow Q\longrightarrow P\longrightarrow1.
\]
These are sequences of full geometric fundamental groups, obtained from the ordered configuration family, as in Part~E; no congruence exactness is assumed. Its geometric monodromy preserves the forgetful fibrations and diagonal inertia, so its image in $\operatorname{Out}(Q)$ is FC-admissible in the sense of Hoshi--Mochizuki. The injectivity assertion of their Theorem~6.1, with all primes and $n=1$, applies to this closed genus-two curve and gives
\[
 \operatorname{Out}^{\rm FC}(Q)\hookrightarrow
 \operatorname{Out}^{\rm FC}(P).
\]
We opened that theorem and its proof, including its reduction to Corollary~5.3; it allows a proper hyperbolic curve at this value of $n$ \cite{D_HoshiMochizuki2011}.

Let $y\in\widehat\Gamma_{2,0}$ act trivially in $\operatorname{Out}(P)$. The displayed injection makes its outer action on $Q$ trivial. Lift $y$ to $z\in\widehat\Gamma_{2,2}$. Conjugation by $z$ on $Q$ is inner, so multiplication by a suitable element of $Q$ makes $z$ centralize $Q$, without changing its image $y$. In particular $z$ centralizes the distinguished subgroup $F$. Its image $z_1\in\widehat\Gamma_{2,1}$ therefore acts trivially in $\operatorname{Out}(F)$: this action is precisely restriction of conjugation by a lift in $\widehat\Gamma_{2,2}$. The already established open genus-two congruence theorem forces $z_1=1$, and its image $y$ is also $1$. This proves the closed case, with the exact groups and the necessary conjugator adjustment specified.

An isomorphism of these two completions means exactly that finite-index congruence subgroups are cofinal among all finite-index subgroups. Restricting the same quotient-complex diagram to a cofinal index system gives the asserted isomorphism of profinite complexes. D.11 then gives the stated connectedness and fundamental groups. At complexity one the unmodified curve complex is disconnected; D.12 and D.21 settle complexity two. This argument asserts no unproved higher profinite connectivity.
\end{proof}

\begin{result}{D.20}{Refuted}{Two literal centralizer assertions in the retrieved Boggi source}
Lemma~6.7 of the retrieved 12 July 2021 version of \emph{Congruence topologies on the mapping class group} asserts that the hyperelliptic involution is self-centralizing in every finite covering-transformation semidirect product specified there. The characteristic mod-two homology cover refutes that assertion. Corollary~4.3 of the same version, when stated for arbitrary unequal multitwist exponents as printed, also cannot identify the centralizer with the full stabilizer of the underlying curve set. Neither assertion is used in D.19.
\end{result}
\begin{proof}
For a closed hyperelliptic surface of genus $g\geq2$, choose a Weierstrass basepoint and set
\[
 N=\ker\bigl(\pi_1(S_g)\longrightarrow H_1(S_g;\mathbb F_2)\bigr).
\]
This subgroup is proper, normal, finite index and characteristic, so it satisfies every subgroup hypothesis in the printed Lemma~6.7 \cite[pp.~26--27]{D_Boggi2020}. The hyperelliptic involution acts as $-I$ on integral homology and hence trivially on the quotient
$G_N=(C_2)^{2g}$. Thus the specified deck group is
\[
 G=G_N\rtimes\langle\iota\rangle
   =(C_2)^{2g+1},\qquad C_G(\iota)=G.
\]
Its centralizer has order $2^{2g+1}$, whereas $\langle\iota\rangle$ has order two. The unramified surface cover has degree $2^{2g}$ and genus $1+2^{2g}(g-1)$ by Euler characteristic. The saved exact computation gives degree $16$, genus $17$, and centralizer order $32$ at $g=2$.

For the second assertion choose disjoint nonseparating curves $a_1,a_2$ on the two handles of a closed genus-two surface, and an orientation-preserving mapping class $s$ exchanging the handles. Then $s$ belongs to the setwise stabilizer of $\{a_1,a_2\}$, which is the stabilizer explicitly defined in Theorem~3.6 of that source. But
\[
 s(T_{a_1}T_{a_2}^{\,2})s^{-1}
     =T_{a_1}^{\,2}T_{a_2}
     \ne T_{a_1}T_{a_2}^{\,2}.
\]
The inequality is visible already on $H_1(S_2;\mathbb F_3)$: in the symplectic basis $(a_1,b_1,a_2,b_2)$ the two matrices have upper-right entries $(1,2)$ and $(2,1)$ in their two diagonal $2$-by-$2$ blocks. Thus it persists in the procongruence group, without any appeal to a centralizer theorem. Both examples were executed in \path{code/d_hyperelliptic_counterexample.py}. The first refutes a universal finite-quotient statement; it does not refute the separate profinite self-normalizer assertion, which requires a cofinal argument. The second shows exactly why a weight-preserving stabilizer must appear when unequal exponents are used.
\end{proof}

\begin{result}{D.21}{Proved}{The second complexity-two profinite obstruction}
The full finite-level completion of the unmodified curve complex $\mathcal C_{1,2}$ has nonzero continuous first cohomology with $\mathbb F_2$ coefficients, and nontrivial profinite fundamental group. By D.19 the same holds for its procongruence completion.
\end{result}
\begin{proof}
Use the torus involution exchanging the two marked points. Its quotient is a sphere with four branch points and a fifth distinguished point, the image of the two marks. Lifting a sphere curve gives an isomorphism
\[
 P^*:\mathcal C_{0,5}\xrightarrow{\sim}\mathcal C_{1,2}.
\]
The precise map and invariant-representative proof are Luo, Lemma~2.1(a), opened on pages~3--4 \cite{D_Luo2000}. A sphere curve separating two branch points lifts to two parallel nonseparating curves and hence one isotopy class. A curve separating the distinguished point and one branch point lifts to one separating curve bounding the disk containing both torus marks. Hyperelliptic invariance supplies an invariant representative of every torus curve; taking invariant disjoint representatives proves the converse disjointness implication. These operations give the displayed inverse correspondence.

The construction is equivariant for the lift of the pure sphere mapping class group. Algebraically, B.9 identifies
\[
 \Gamma_{1,2}=B_4/\langle\Delta_4^2\rangle,\qquad
 \Gamma_{0,5}=P_4/\langle\Delta_4^2\rangle,
\]
and the latter subgroup has index $24$, by the permutation homomorphism to $\mathfrak S_4$. The subgroup is exactly the one retaining all four branch labels in the preceding cover construction. Finite-index subgroups of this subgroup are cofinal in the finite-index system for $\Gamma_{1,2}$; normality can be imposed by taking the finite intersection of conjugates. The equivariant simplicial isomorphism therefore identifies the two full finite-level curve complexes. Transporting the explicit continuous Petersen-edge class and the closed pentagon of D.12 gives a nonzero continuous $\mathbb F_2$ class on the genus-one complex. Its profinite fundamental group is consequently nontrivial. D.19 supplies the final procongruence comparison.
\end{proof}

